%% ============================================================
%% Cartographie géotechnique nationale par méthodes géostatistiques multi-modèles
%% Application au Togo — KED Hiérarchique, RK-SCORPAN, Fusion BLUP, VfS, MTGP
%% ============================================================
\documentclass[10pt,a4paper,twocolumn]{article}

\usepackage[a4paper,top=2.2cm,bottom=2.2cm,left=1.5cm,right=1.5cm,columnsep=0.8cm]{geometry}
\usepackage[french]{babel}
\usepackage{csquotes}
\usepackage{graphicx}
\usepackage{amsmath,amssymb,amsfonts}
\usepackage{bm}
\usepackage{microtype}
\usepackage{titlesec}
\usepackage{caption}
\usepackage{subcaption}
\usepackage{booktabs}
\usepackage{array}
\usepackage{multirow}
\usepackage{placeins}
\usepackage{afterpage}
\usepackage{xcolor}
\usepackage{hyperref}
\hypersetup{hidelinks}
\usepackage{fontspec}
\usepackage{newtxmath}
\usepackage[style=ieee,backend=biber]{biblatex}
\addbibresource{bibliographie.bib}

\IfFontExistsTF{Times New Roman}{
  \setmainfont{Times New Roman}
}{\setmainfont{TeX Gyre Termes}}

%% Mise en page sections
\titleformat{\section}{\normalfont\fontsize{11}{13}\bfseries}{\thesection.}{0.8em}{}
\titleformat{\subsection}{\normalfont\fontsize{10}{12}\bfseries}{\thesubsection}{0.8em}{}
\titleformat{\subsubsection}{\normalfont\fontsize{10}{12}\itshape}{\thesubsubsection}{0.6em}{}
\titlespacing{\section}{0pt}{8pt}{4pt}
\titlespacing{\subsection}{0pt}{6pt}{3pt}

%% Légendes
\captionsetup{font=small, labelfont=bf, labelsep=endash, justification=justified}

%% Commandes mathématiques pratiques
\newcommand{\bZ}{\mathbf{Z}}
\newcommand{\bx}{\mathbf{x}}
\newcommand{\bs}{\mathbf{s}}
\newcommand{\bK}{\mathbf{K}}
\newcommand{\blambda}{\boldsymbol{\lambda}}
\newcommand{\bSigma}{\boldsymbol{\Sigma}}
\newcommand{\KED}{\text{KED-H}}
\newcommand{\RK}{\text{RK-SCORPAN}}
\newcommand{\BLUP}{\text{BLUP}}

%% En-tête du titre
\title{%
  \vspace{-1.4cm}%
  \fontsize{13}{15}\bfseries%
  CARTOGRAPHIE GÉOTECHNIQUE NATIONALE PAR MÉTHODES\\
  GÉOSTATISTIQUES MULTI-MODÈLES : KRIGEAGE AVEC DÉRIVE\\
  HIÉRARCHIQUE, RÉGRESSION SCORPAN ET FUSION BAYÉSIENNE\\[4pt]
  {\fontsize{11}{13}\normalfont\itshape Application au Territoire Togolais
    (56\,600 km², 122 sondages, 5 paramètres)}%
}

\author{%
  \fontsize{10}{12}\normalfont
  Serge TABE DJATO\textsuperscript{1*}\\[3pt]
  \fontsize{9}{11}\normalfont
  \textsuperscript{1}Institut des Sciences Technologiques Économiques et Administratives
  (FORMATEC), Lomé, Togo.\\
  \textsuperscript{*}Correspondance : \texttt{serge.tdjato@gmail.com}
}
\date{}

\begin{document}

\twocolumn[%
  \maketitle
  \vspace{-0.5cm}
  \rule{\linewidth}{0.4pt}
  \vspace{0.15cm}

  % ── RÉSUMÉ ────────────────────────────────────────────────────────────────
  \noindent\small
  \textbf{Résumé ---} \textit{%
    La cartographie géotechnique à l'échelle nationale pose un problème fondamental
    de densité : 122 sondages pour 56\,600~km\textsuperscript{2} de territoire togolais.
    Cet article présente une hiérarchie de modèles géostatistiques pour interpoler
    cinq paramètres géotechniques (VBS, IP, WL, WP, EG) sur 29\,407 mailles
    (2~km $\times$ 2~km). Le Krigeage avec Dérive Hiérarchique (KED-H) intègre cinq
    niveaux de connaissance géologique, réduisant le LOO-RMSE du VBS à 2{,}941~g/100g.
    La Régression Kriging SCORPAN (RK) atteint 2{,}480~g/100g. Une Fusion Bayésienne
    BLUP réduit la variance de 46{,}7\,\%\,$\pm$\,1{,}8\,\%. Un modèle PLS sur indices
    spectraux Sentinel-2 (VfS) et un Processus Gaussien Multi-Tâches
    ($r_{\mathrm{IP,WL}}{=}0{,}802$) complètent la hiérarchie.
    L'architecture produit 441\,105 valeurs interpolées avec intervalles de confiance.%
  }

  \smallskip
  \noindent\small
  \textbf{Abstract ---} \textit{%
    National-scale geotechnical mapping presents a fundamental information-density
    challenge: 122 boreholes over 56\,600~km\textsuperscript{2} of Togolese territory.
    This paper presents a hierarchy of geostatistical models interpolating five
    geotechnical parameters (VBS, IP, WL, WP, EG) over 29,407 cells (2~km $\times$ 2~km).
    Hierarchical KED (KED-H) reduces the VBS LOO-RMSE to 2.941~g/100g; SCORPAN
    Regression Kriging (RK) achieves 2.480~g/100g. Bayesian BLUP fusion reduces
    variance by 46.7\,\%\,$\pm$\,1.8\,\%. Sentinel-2 PLS and Multi-Task Gaussian
    Process ($r_{\mathrm{IP,EG}}{=}0.742$) complete the hierarchy.
    The system produces 441,105 interpolated values with confidence intervals.%
  }

  \smallskip
  \noindent\small
  \textbf{Mots clés :} Krigeage Externe, SCORPAN, Fusion Bayésienne, MTGP,
    Géotechnique, Togo.
  \quad\textbf{Keywords:} External Drift Kriging, SCORPAN, Bayesian Fusion, MTGP,
    Geotechnics, Togo.

  \vspace{0.15cm}
  \rule{\linewidth}{0.4pt}
  \vspace{0.3cm}
]

% ══════════════════════════════════════════════════════════════════════════════
\section{Introduction}
% ══════════════════════════════════════════════════════════════════════════════

La planification des infrastructures civiles en Afrique subsaharienne repose
quasi-exclusivement sur des données géotechniques ponctuelles, dont la densité
spatiale est incompatible avec les exigences du dimensionnement distribué.
Au Togo, les cinq paramètres géotechniques fondamentaux --- la Valeur au Bleu
de Méthylène (VBS), l'Indice de Plasticité (IP), la Limite de Liquidité (WL),
la Limite de Plasticité (WP) et le Coefficient de Gonflement (EG) --- ne sont
disponibles qu'en 122 points de mesure pour un territoire de 56\,600~km², soit
un sondage pour 463~km² \cite{mcbratney2003digital}.

La transformation de cette information lacunaire en une surface de prédiction
continue et quantifiée est un problème inverse mal posé, classiquement résolu
par interpolation géostatistique. Les travaux fondateurs de Matheron
\cite{matheron1963} ont établi le Krigeage Ordinaire (OK) comme estimateur
sans biais de variance minimale. Des extensions successives ont incorporé
des variables auxiliaires : le Krigeage avec Dérive Externe (KED)
\cite{odeh1994further} utilise une variable explicative continue comme composante
de dérive non-stationnaire, tandis que la Régression Kriging (RK)
\cite{hengl2007regression} sépare explicitement la tendance déterministe,
estimée par régression, du résidu stochastique estimé par krigeage ordinaire.

Dans le contexte pédogéologique togolais, ces méthodes présentent deux limites
structurelles. D'une part, les couches géologiques disponibles (unités pédologiques,
formations géologiques, risque de gonflement RGA) offrent une information
catégorielle multi-échelle que ni KED ni RK n'exploitent pleinement. D'autre part,
les cinq paramètres géotechniques sont fortement corrélés entre eux
($r_{\mathrm{IP,WL}} = 0{,}802$, $r_{\mathrm{IP,EG}} = 0{,}742$) : leur modélisation
indépendante ignore une information inter-paramètre substantielle.

Cet article présente une hiérarchie de quatre méthodes développée pour répondre
à ces limites : (L1) le Krigeage avec Dérive Hiérarchique à cinq niveaux (KED-H),
(L2a) la Régression Kriging SCORPAN (RK-SCORPAN), (L2b) leur Fusion Bayésienne
par BLUP, (L3) la Régression PLS sur indices spectraux Sentinel-2 (VfS-PLS), et
(L4) le Processus Gaussien Multi-Tâches par Modèle de Coregionalisation Intrinsèque
(MTGP/ICM). L'ensemble est validé par validation croisée Leave-One-Out (LOO-CV)
sur les cinq paramètres et les trois horizons de profondeur.

La section~\ref{sec:donnees} décrit les données et le contexte géologique.
La section~\ref{sec:modeles} développe les formulations mathématiques complètes.
La section~\ref{sec:pipeline} présente l'architecture du système d'information.
La section~\ref{sec:sensibilite} analyse la sensibilité et la robustesse.
La section~\ref{sec:covariables} détaille le calcul des covariables SCORPAN.
La section~\ref{sec:algo} expose l'implémentation algorithmique.
La section~\ref{sec:resultats} présente les résultats et la validation.
La section~\ref{sec:geologie} analyse l'interprétation géologique.
La section~\ref{sec:perspectives} discute les développements futurs.
La section~\ref{sec:conclusion} conclut et synthétise les contributions.

% ══════════════════════════════════════════════════════════════════════════════
\section{Rappels théoriques sur la géostatistique linéaire}
\label{sec:theorie}
% ══════════════════════════════════════════════════════════════════════════════

\subsection{Hypothèses de stationnarité}

Le cadre géostatistique repose sur des hypothèses de stationnarité qui
conditionnent la validité des estimateurs. La stationnarité d'ordre 2
(intrinsèque dans sa formulation la plus faible) requiert que :

\begin{equation}
  \mathbb{E}[Z(\bs)] = m(\bs), \quad
  \mathrm{Cov}[Z(\bs), Z(\bs')] = C(\bs - \bs')
  \label{eq:stationnarite}
\end{equation}

c'est-à-dire que l'espérance peut être non-constante (dérive $m(\bs)$) mais
que la covariance ne dépend que de la distance vectorielle $\bs - \bs'$.
Dans les données géotechniques togolaises, cette hypothèse est approximativement
satisfaite après déduction de la dérive hiérarchique $\mu_h(\bs)$
\eqref{eq:prior_hierarchique}.

La stationnarité intrinsèque, formulée en termes de semi-variogramme :

\begin{equation}
  \gamma(\mathbf{h}) = \frac{1}{2} \mathrm{Var}[Z(\bs + \mathbf{h}) - Z(\bs)]
  \label{eq:intrinsec}
\end{equation}

est moins restrictive et couvre les cas où la covariance $C(h)$ n'existe pas
formellement (processus à variance infinie), tel que le mouvement brownien
fractionnaire. En pratique, les données géotechniques satisfont la stationnarité
d'ordre 2 après déduction de la dérive, et les deux formulations sont équivalentes.

\subsection{Propriétés d'optimalité du krigeage}

L'estimateur de krigeage $Z^* = \sum_i \lambda_i z(\bs_i)$ est optimal au sens
de la variance quadratique minimale (MVUE pour \textit{Minimum Variance Unbiased
Estimator}) dans la classe des estimateurs linéaires non-biaisés. Ces trois propriétés
sont simultanément garanties par les équations de krigeage~\eqref{eq:ked_system} :

\begin{enumerate}
  \item \textbf{Linéarité} : $Z^*(\bs_0)$ est une combinaison linéaire des
    observations $z(\bs_i)$.

  \item \textbf{Non-biais} : $\mathbb{E}[Z^*(\bs_0) - Z(\bs_0)] = 0$,
    garanti par les contraintes de Lagrange sur les multiplicateurs $\bm{\mu}$.

  \item \textbf{Variance minimale} : $\mathrm{Var}[Z^*(\bs_0) - Z(\bs_0)]$
    est minimale parmi tous les estimateurs linéaires non-biaisés.
\end{enumerate}

Cette optimalité est conditionnelle au modèle de covariance $C(h)$ choisi.
Si le modèle est incorrect, les poids $\lambda_i$ sont sous-optimaux et la
variance de krigeage~\eqref{eq:ked_var} sous-estime l'erreur réelle, ce qui
motive la robustesse aux choix variographiques analysée en section~\ref{sec:sensibilite}.

\subsection{Relation entre krigeage et régression}

Il est instructif d'établir la relation entre le krigeage et la régression
linéaire multiple. Dans le cas limite où la portée $a \to 0$ (pépite pure,
$\gamma(h) = C_0$ pour $h > 0$), le krigeage ordinaire dégénère en la
moyenne arithmétique des observations. À l'inverse, quand $a \to \infty$,
la matrice de covariance $\mathbf{K}$ devient $C_0 \mathbf{1}\mathbf{1}^T$
et le krigeage ordinaire converge vers la régression de la variable $z$
sur une constante (la moyenne des observations).

Le KED avec dérive polynomiale d'ordre 1 est formellement équivalent à la
Régression Kriging dans la limite où l'espérance de krigeage est modélisée
comme une tendance linéaire. La distinction entre les deux approches réside
dans le traitement de l'inbiaisabilité : le KED impose les contraintes dans
l'espace primal, la RK les impose via la décomposition de la prédiction en
tendance déterministe et résidu aléatoire \cite{hengl2007regression}.

\subsection{Intervalles de prédiction vs intervalles de confiance}

La variance de krigeage $\sigma^2_K(\bs_0)$ fournit un \emph{intervalle
de prédiction} pour une réalisation individuelle $Z(\bs_0)$, non un intervalle
de confiance sur la moyenne. En supposant la gaussianité des erreurs :

\begin{equation}
  P\bigl[Z(\bs_0) \in Z^*(\bs_0) \pm z_{0,975}\, \sigma_K(\bs_0)\bigr] \approx 95\%
  \label{eq:intervalle_pred}
\end{equation}

où $z_{0,975} = 1{,}96$. Cette interprétation est celle adoptée dans les
figures de transect (Fig.~\ref{fig:transect}). La gaussianité est approximativement
vérifiée pour WL, WP et IP, mais non pour VBS (distribution asymétrique à droite) ;
pour ce paramètre, les intervalles de prédiction sont légèrement libéraux
et une transformation logarithmique améliorerait la calibration.

% ══════════════════════════════════════════════════════════════════════════════
\section{Zone d'étude, données et contexte géologique}
\label{sec:donnees}
% ══════════════════════════════════════════════════════════════════════════════

\subsection{Territoire togolais et contraintes spatiales}

Le Togo s'étend sur 56\,600~km\textsuperscript{2} entre 6°N et 11°N,
avec une largeur maximale de 145~km. Cinq grandes régions physiographiques
structurent la variabilité géotechnique : la zone maritime (alluvions
quaternaires côtières), les plateaux (ferralisols sur socle précambrien),
la région centrale (gneiss et orthogneiss, altitude 800–900~m), la région
de Kara (schistes et quartzites de l'Atacora), et les savanes septentrionales
(grès voltaïens, sub-sahélien). Cette structure impose une non-stationnarité
spatiale marquée des paramètres géotechniques, validée par le semi-variogramme
empirique omnidirectionnel du VBS (Fig.~\ref{fig:variogram}).

La grille nationale d'interpolation est composée de $N_g = 29\,407$ mailles
carrées de 2~km~×~2~km couvrant l'intégralité du territoire. La densité
de sondages est $\rho = 122/56\,600 \approx 2{,}2 \times 10^{-3}$~sondage~km$^{-2}$,
valeur caractéristique des pays en développement pour lesquels la cartographie
géotechnique numérique constitue un enjeu majeur.

\subsection{Base de données géotechniques}

La base de données comprend 122 sondages géolocalisés ($\pm$2~m GPS),
distribués sur l'ensemble du territoire avec une densité plus élevée
dans les régions maritime et centrale. Pour chaque sondage, jusqu'à trois
horizons de profondeur sont documentés :

\begin{itemize}
  \item H1 : 0–1~m (sol superficiel, couche de fondation habituelle)
  \item H2 : 1–1,5~m (transition vers le sol résiduel)
  \item H3 : 1,5–2~m (sol résiduel profond)
\end{itemize}

Le tableau~\ref{tab:stats_obs} présente les statistiques descriptives
des cinq paramètres mesurés. La distribution du VBS est fortement asymétrique
(médiane = 3,33~g/100g, maximum = 28,86~g/100g), reflétant la présence de
Vertisols argileux gonflants dans certaines zones. L'IP montre une
variabilité modérée ($\sigma = 13{,}2\%$) cohérente avec la diversité
minéralogique des formations traversées.

\begin{table}[!ht]
\centering\small
\caption{Statistiques descriptives des paramètres géotechniques mesurés in situ
  ($n_{\mathrm{VBS}} = 329$, $n_{\mathrm{WL/WP}} \approx 360$,
  $n_{\mathrm{EG}} = 327$).}
\label{tab:stats_obs}
\begin{tabular}{lrrrrrr}
\toprule
Param. & $\bar{x}$ & $\sigma$ & $Q_1$ & Méd. & $Q_3$ & Max \\
\midrule
VBS (g/100g) & 4,27 & 3,91 & 1,64 & 3,33 & 5,71 & 28,86 \\
IP (\%)       & 21,5 & 13,2 & 11,0 & 19,8 & 29,6 & 68,0  \\
WL (\%)       & 40,9 & 12,6 & 31,5 & 39,8 & 48,7 & 72,9  \\
WP (\%)       & 19,5 & 7,9  & 14,2 & 18,5 & 23,5 & 42,0  \\
EG (\%)       & 4,30 & 1,82 & 3,10 & 4,20 & 5,50 & 10,2  \\
\bottomrule
\end{tabular}
\end{table}

\begin{figure}[!ht]
  \centering
  \includegraphics[width=\linewidth]{figures/fig01_boxplots_parametres.pdf}
  \caption{Distribution des paramètres géotechniques mesurés in situ.
    La ligne centrale représente la médiane, la boîte l'intervalle
    inter-quartile [Q1, Q3], et les points les valeurs aberrantes.}
  \label{fig:boxplots}
\end{figure}

\subsection{Couches géologiques et pédologiques auxiliaires}
\label{subsec:couches}

Cinq couches spatiales constituent la base des covariables géologiques,
disponibles à couverture nationale complète :

\begin{enumerate}
  \item \textbf{Unités géologiques} : 120 polygones, 15 formations
    (gneiss socle, orthogneiss, schistes, quartzites, alluvions, cuirasses) ;
  \item \textbf{Unités pédologiques} : 61 polygones, 13 types de sol
    (Vertisols, Sols ferralitiques, Sols ferrugineux, Sols hydromorphes…) ;
  \item \textbf{Risque de gonflement RGA} \cite{chassagneux1996} :
    342 entités, 5 niveaux de risque (Très faible à Très élevé) ;
  \item \textbf{Hydrogéologie} : 14 polygones, profondeur et perméabilité
    de la nappe phréatique ;
  \item \textbf{Zones géomorphologiques spéciales} : 5 zones identifiées
    (Dépression de la Lama, Bado, Plaine du Mono, Oti, Fosse aux Lions).
\end{enumerate}

La pertinence géotechnique de ces couches repose sur la minéralogie argileuse :
les gneiss et micaschistes (53\% du territoire) produisent par altération
des kaolinites (VBS modéré, EG faible), tandis que les argiles sédimentaires
et les shales génèrent des smectites gonflantes (VBS élevé, EG fort). Cette
relation est quantifiée via le VBS moyen par type de sol \cite{chassagneux1996}.

% ══════════════════════════════════════════════════════════════════════════════
\section{Formulations mathématiques des modèles}
\label{sec:modeles}
% ══════════════════════════════════════════════════════════════════════════════

Soit $Z(\bs)$ le champ aléatoire d'un paramètre géotechnique en tout point
$\bs \in \mathbb{R}^2$ du territoire. On dispose d'observations
$\{z(\bs_i)\}_{i=1}^{n}$ en $n$ sondages géolocalisés. L'objectif est
de prédire $Z(\bs_0)$ en tout point $\bs_0$ de la grille nationale,
avec quantification de l'incertitude $\sigma^2(\bs_0) = \mathrm{Var}[Z^*(\bs_0) - Z(\bs_0)]$.

\subsection{L1 — Krigeage avec Dérive Externe Hiérarchique (KED-H)}
\label{subsec:ked}

\subsubsection{Modèle structural}

Le KED-H décompose le champ $Z(\bs)$ en une composante de dérive
non-stationnaire $m(\bs)$ et un résidu stationnaire $R(\bs)$ de moyenne nulle :

\begin{equation}
  Z(\bs) = m(\bs) + R(\bs), \quad \mathbb{E}[R(\bs)] = 0
  \label{eq:ked_decomp}
\end{equation}

La dérive est modélisée comme une combinaison linéaire de fonctions connues
$\{f_k(\bs)\}_{k=0}^{K}$, ici les moyennes par contexte géologique hiérarchique :

\begin{equation}
  m(\bs) = \sum_{k=0}^{K} a_k f_k(\bs)
  \label{eq:ked_drift}
\end{equation}

\subsubsection{Prior hiérarchique à cinq niveaux}

La variable de dérive $f(\bs)$ est construite par repli hiérarchique sur
cinq niveaux de connaissance géologique, du plus spécifique au plus général :

\begin{equation}
  \mu_h(\bs) = \begin{cases}
    \bar{z}_{\mathrm{zone}} & \text{si } \bs \in \mathcal{Z}_k,\, n_k \geq n_{\min} \\
    \bar{z}_{\mathrm{ped}}  & \text{si type pédologique disponible} \\
    \bar{z}_{\mathrm{RGA}}  & \text{si risque Chassagneux disponible} \\
    \bar{z}_{\mathrm{géol}} & \text{si formation géologique disponible} \\
    \bar{z}_{\mathrm{nat}}  & \text{moyenne nationale (repli ultime)}
  \end{cases}
  \label{eq:prior_hierarchique}
\end{equation}

avec $n_{\min} = 5$ sondages par catégorie requis pour estimer la moyenne
de manière robuste. Cette hiérarchie garantit que la connaissance la plus
fine disponible est utilisée, avec repli automatique en cas de données
insuffisantes.

\subsubsection{Équations de krigeage}

L'estimateur KED est :

\begin{equation}
  Z^*_{\KED}(\bs_0) = \sum_{i=1}^{n} \lambda_i(\bs_0)\, z(\bs_i)
  \label{eq:ked_estimateur}
\end{equation}

Les poids $\blambda = (\lambda_1, \ldots, \lambda_n)^T$ sont solution du
système linéaire obtenu par les conditions d'inbiaisabilité et de
minimisation de la variance :

\begin{equation}
  \begin{pmatrix} \bK & \mathbf{F} \\ \mathbf{F}^T & \mathbf{0} \end{pmatrix}
  \begin{pmatrix} \blambda \\ \bm{\mu} \end{pmatrix} =
  \begin{pmatrix} \mathbf{k}_0 \\ \mathbf{f}_0 \end{pmatrix}
  \label{eq:ked_system}
\end{equation}

où $\bK_{ij} = C(\bs_i, \bs_j)$ est la matrice de covariance des données,
$\mathbf{k}_0 = (C(\bs_i, \bs_0))_{i=1}^n$ le vecteur de covariance
données-cible, $\mathbf{F}_{ik} = f_k(\bs_i)$ la matrice des fonctions
de dérive aux points de données, $\mathbf{f}_0 = (f_k(\bs_0))_{k=0}^K$
les fonctions de dérive au point cible, et $\bm{\mu}$ les multiplicateurs
de Lagrange assurant les contraintes d'inbiaisabilité.

La variance de krigeage associée s'écrit :

\begin{equation}
  \sigma^2_{\KED}(\bs_0) = C(0) - \begin{pmatrix} \blambda \\ \bm{\mu} \end{pmatrix}^T
  \begin{pmatrix} \mathbf{k}_0 \\ \mathbf{f}_0 \end{pmatrix}
  \label{eq:ked_var}
\end{equation}

\subsubsection{Modèle variographique}

La structure spatiale du résidu $R(\bs)$ est caractérisée par le
semi-variogramme $\gamma(h) = \frac{1}{2}\mathrm{E}[(R(\bs+h) - R(\bs))^2]$.
Un modèle sphérique est ajusté aux données empiriques du VBS H1 :

\begin{equation}
  \gamma(h) = C_0 + C \cdot \begin{cases}
    \dfrac{3h}{2a} - \dfrac{1}{2}\!\left(\dfrac{h}{a}\right)^{\!3} & h \leq a \\[6pt]
    1 & h > a
  \end{cases}
  \label{eq:variogram_spherique}
\end{equation}

avec les paramètres calibrés : effet de pépite $C_0 = 2{,}8$~(g/100g)²,
palier $C_0 + C = 14{,}0$~(g/100g)², portée $a = 220$~km.
La figure~\ref{fig:variogram} illustre l'ajustement du modèle aux
semi-variances empiriques ainsi que l'anisotropie géométrique observée.

\begin{figure}[!ht]
  \centering
  \includegraphics[width=\linewidth]{figures/fig02_variogram_vbs_h1.pdf}
  \caption{Analyse variographique du VBS à l'horizon H1. \textit{Gauche :}
    semi-variogramme omnidirectionnel empirique (points, taille proportionnelle
    au nombre de paires) et modèle sphérique ajusté. \textit{Droite :}
    anisotropie directionnelle N-S / E-O.}
  \label{fig:variogram}
\end{figure}

\subsection{L2a — Régression Kriging SCORPAN (RK)}
\label{subsec:rk}

\subsubsection{Cadre conceptuel SCORPAN}

Le modèle SCORPAN \cite{mcbratney2003digital} décompose la propriété du
sol $z$ comme fonction de six facteurs de formation :

\begin{equation}
  z(\bs) = f\bigl(S, C, O, R, P, A, N\bigr) + \varepsilon(\bs)
  \label{eq:scorpan}
\end{equation}

où $S$ = propriétés intrinsèques du sol, $C$ = climat, $O$ = organismes,
$R$ = relief (DSM, pente, TPI, HAND), $P$ = matériau parental (géologie),
$A$ = âge, $N$ = position spatiale (coordonnées). Le terme résiduel
$\varepsilon(\bs)$ est spatialement corrélé.

\subsubsection{Régression par Ridge et covariables}

La composante déterministe est estimée par régression Ridge pénalisée,
robuste à la colinéarité des covariables terrain :

\begin{equation}
  \hat{\bm{\beta}} = \arg\min_{\bm{\beta}}
  \Bigl\| \mathbf{z} - \mathbf{X}\bm{\beta} \Bigr\|^2 + \lambda_R \|\bm{\beta}\|^2
  \label{eq:ridge}
\end{equation}

où $\mathbf{X} \in \mathbb{R}^{n \times p}$ est la matrice des covariables
et $\lambda_R$ le paramètre de régularisation calibré par validation croisée.
La solution fermée est :

\begin{equation}
  \hat{\bm{\beta}} = (\mathbf{X}^T\mathbf{X} + \lambda_R \mathbf{I})^{-1} \mathbf{X}^T \mathbf{z}
  \label{eq:ridge_sol}
\end{equation}

Les $p = 12$ covariables utilisées incluent : altitude DSM, pente, indice
TPI (relief relatif), HAND (hauteur au-dessus du réseau hydrographique),
précipitations annuelles et saisonnières (WorldClim), formation géologique
(encodage ordinal), type pédologique (encodage), risque RGA, coordonnées
UTM31N ($x$, $y$), et indice NDVI médian Sentinel-2.

\subsubsection{Krigeage des résidus}

Les résidus de régression aux points d'entraînement :

\begin{equation}
  \hat{r}(\bs_i) = z(\bs_i) - \mathbf{x}(\bs_i)^T \hat{\bm{\beta}}, \quad i = 1, \ldots, n
  \label{eq:residus_rk}
\end{equation}

sont interpolés par krigeage ordinaire. L'estimateur RK complet est :

\begin{equation}
  Z^*_{\RK}(\bs_0) = \mathbf{x}(\bs_0)^T \hat{\bm{\beta}}
    + \sum_{i=1}^{n} \lambda_i^{RK}(\bs_0)\, \hat{r}(\bs_i)
  \label{eq:rk_estimateur}
\end{equation}

La variance de prédiction RK, qui intègre à la fois l'incertitude
de régression et l'incertitude de krigeage des résidus, est :

\begin{equation}
  \sigma^2_{\RK}(\bs_0) = \sigma^2_{r}(\bs_0) + \mathbf{x}(\bs_0)^T
  (\mathbf{X}^T\mathbf{X})^{-1} \mathbf{x}(\bs_0) \hat{\sigma}^2_\varepsilon
  \label{eq:rk_var}
\end{equation}

où $\sigma^2_{r}(\bs_0)$ est la variance de krigeage des résidus et
$\hat{\sigma}^2_\varepsilon$ la variance résiduelle de régression.

\subsection{L2b — Fusion Bayésienne par BLUP}
\label{subsec:fusion}

\subsubsection{Motivation théorique}

KED et RK capturent des structures complémentaires :
KED exploite la continuité spatiale des résidus et le prior géologique ;
RK capte les tendances déterministes longue portée (topographie, climat).
Leur combinaison optimale par le Meilleur Estimateur Linéaire Sans Biais
(BLUP) est mathématiquement justifiée lorsque les deux estimateurs sont
non-biaisés \cite{chilesdelfiner2012}.

\subsubsection{Combinaison pondérée par incertitude}

Pour tout point $\bs_0$, la Fusion BLUP est définie par :

\begin{equation}
  Z^*_{\BLUP}(\bs_0) =
    \frac{w_{\KED}(\bs_0)\, Z^*_{\KED}(\bs_0) + w_{\RK}(\bs_0)\, Z^*_{\RK}(\bs_0)}
         {w_{\KED}(\bs_0) + w_{\RK}(\bs_0)}
  \label{eq:blup}
\end{equation}

avec les poids inversement proportionnels à la variance de prédiction locale :

\begin{equation}
  w_{\KED}(\bs_0) = \frac{1}{\sigma^2_{\KED}(\bs_0)}, \qquad
  w_{\RK}(\bs_0)  = \frac{1}{\sigma^2_{\RK}(\bs_0)}
  \label{eq:blup_poids}
\end{equation}

Cette formulation équivaut à introduire un coefficient de mélange
$\beta(\bs_0) \in [0,1]$ adaptatif :

\begin{equation}
  \beta(\bs_0) = \frac{\sigma^2_{\KED}(\bs_0)}{\sigma^2_{\KED}(\bs_0) + \sigma^2_{\RK}(\bs_0)}
  \label{eq:beta}
\end{equation}

tel que $Z^*_{\BLUP}(\bs_0) = Z^*_{\KED}(\bs_0) + \beta(\bs_0)\bigl[Z^*_{\RK}(\bs_0) - Z^*_{\KED}(\bs_0)\bigr]$.

La variance de la fusion, toujours inférieure au minimum des deux variances
individuelles, est :

\begin{equation}
  \sigma^2_{\BLUP}(\bs_0) = \frac{1}{w_{\KED}(\bs_0) + w_{\RK}(\bs_0)}
  = \frac{\sigma^2_{\KED} \cdot \sigma^2_{\RK}}{\sigma^2_{\KED} + \sigma^2_{\RK}}
  \label{eq:blup_var}
\end{equation}

Cette propriété garantit que $\sigma^2_{\BLUP} < \min(\sigma^2_{\KED}, \sigma^2_{\RK})$
pour tout $\bs_0$, soit une réduction d'incertitude systématique et spatialement
adaptative \cite{rasmussen2006gaussian}.

\subsection{L3 — Régression PLS sur Indices Spectraux Sentinel-2 (VfS-PLS)}
\label{subsec:vfs}

\subsubsection{Fondements scientifiques}

Les argiles gonflantes (smectites, illites) présentent des absorptions
caractéristiques dans le domaine infrarouge à ondes courtes (SWIR) :
2~200~nm (Al-OH), 2~300~nm (Mg-OH) \cite{viscarra2006data}. Les bandes
spectrales B11 (1\,610~nm) et B12 (2\,190~nm) de Sentinel-2 permettent
de calculer un indice d'activité argileuse :

\begin{equation}
  I_{\mathrm{argile}} = \frac{B_{11}}{B_{12}}
  \label{eq:clay_index}
\end{equation}

La valeur $I_{\mathrm{argile}} > 1$ indique une absorption en B12 plus
forte, caractéristique des argiles gonflantes (smectites). Des indices
complémentaires sont calculés :

\begin{equation}
  I_{\mathrm{SWIR}} = \frac{B_{11} - B_{12}}{B_{11} + B_{12}}, \quad
  \mathrm{NDVI} = \frac{B_8 - B_4}{B_8 + B_4}, \quad
  I_{\mathrm{Fe}} = \frac{B_4}{B_8}
  \label{eq:indices_spectraux}
\end{equation}

\subsubsection{Modèle PLS}

La Régression par Moindres Carrés Partiels (PLS) est la méthode
de référence pour la prédiction de propriétés du sol à partir de
spectres \cite{viscarra2006data}. Pour $n$ sondages avec $p$
variables spectrales ($p \ll n$), le modèle PLS cherche
$q \leq p$ composantes latentes $\mathbf{T} = \mathbf{X}\mathbf{W}^*$
maximisant la covariance entre $\mathbf{X}$ (matrices d'indices spectraux)
et $\mathbf{y}$ (vecteur VBS) :

\begin{equation}
  \mathbf{w}_k^* = \arg\max_{\|\mathbf{w}\|=1} \, \mathrm{Cov}^2(\mathbf{X}_{k-1}\mathbf{w}, \mathbf{y}_{k-1})
  \label{eq:pls_direction}
\end{equation}

L'estimateur PLS est :

\begin{equation}
  \hat{y}_{\mathrm{PLS}}(\bs_0) = \bar{y} + \mathbf{x}(\bs_0)^T \mathbf{W}^*(\mathbf{P}^T\mathbf{W}^*)^{-1}\mathbf{q}
  \label{eq:pls_pred}
\end{equation}

où $\mathbf{W}^*$ est la matrice des vecteurs de poids, $\mathbf{P}$ les
chargements X, et $\mathbf{q}$ les chargements Y. Le nombre de composantes
($q = 3$ retenu) est optimisé par LOO-CV.

\subsection{L4 — Processus Gaussien Multi-Tâches (MTGP/ICM)}
\label{subsec:mtgp}

\subsubsection{Modèle de Coregionalisation Intrinsèque}

Les cinq paramètres géotechniques $(Z_1, \ldots, Z_5)$ = (VBS, IP, WL, WP, EG)
présentent des corrélations structurelles fortes (tableau~\ref{tab:correlations}).
Le Modèle de Coregionalisation Intrinsèque (ICM) \cite{journelhijbregts1978}
représente chaque paramètre comme combinaison linéaire de $r$ processus
gaussiens latents indépendants $\{G_l(\bs)\}_{l=1}^r$ :

\begin{equation}
  Z_d(\bs) = \sum_{l=1}^{r} a_{dl}\, G_l(\bs) + \varepsilon_d(\bs), \quad d = 1,\ldots,5
  \label{eq:icm}
\end{equation}

En notation matricielle, la covariance conjointe inter-paramètres est :

\begin{equation}
  \mathrm{Cov}[Z_d(\bs), Z_{d'}(\bs')] = \sum_{l=1}^{r} a_{dl} a_{d'l}\, k_l(\bs, \bs')
  \label{eq:icm_cov}
\end{equation}

où $k_l(\bs, \bs') = \sigma_l^2 \exp(-\|\bs - \bs'\|/\ell_l)$ est
le noyau de base Matérn 3/2 adapté aux discontinuités géologiques.
La matrice $\mathbf{A} = (a_{dl})$ est la décomposition de Cholesky
de la matrice de coregionalisation, garantissant la semi-définie positivité
de la covariance conjointe.

\subsubsection{Avantage du co-krigeage}

L'intérêt du MTGP est d'emprunter la structure spatiale d'un paramètre
bien documenté pour améliorer la prédiction d'un autre moins échantillonné.
Avec $n_{\mathrm{IP}} = 93$ et $n_{\mathrm{EG}} = 64$ mesures H1, la
corrélation $r_{\mathrm{IP,EG}} = 0{,}742$ permet à IP d'informer EG.
La prédiction conjointe s'effectue par résolution du système linéaire
généralisé :

\begin{equation}
  \begin{pmatrix} \mathbf{K}_{11} & \cdots & \mathbf{K}_{15} \\
                  \vdots          & \ddots & \vdots          \\
                  \mathbf{K}_{51} & \cdots & \mathbf{K}_{55} \end{pmatrix}
  \begin{pmatrix} \blambda_1 \\ \vdots \\ \blambda_5 \end{pmatrix} =
  \begin{pmatrix} \mathbf{k}_{01} \\ \vdots \\ \mathbf{k}_{05} \end{pmatrix}
  \label{eq:mtgp_system}
\end{equation}

où $\mathbf{K}_{dd'} = \mathrm{Cov}[Z_d(\bs_i), Z_{d'}(\bs_j)]$ selon
l'équation~\eqref{eq:icm_cov}. La complexité algorithmique de l'inversion
est $\mathcal{O}((nD)^3)$ où $D=5$, soit $\mathcal{O}(n^3)$ en pratique
pour les tailles d'entraînement considérées ($n \leq 600$), rendant
le calcul exact (non variationnel) \cite{rasmussen2006gaussian}.

\begin{table}[!ht]
\centering\small
\caption{Matrice de corrélation de Pearson des paramètres géotechniques
  (mesures in situ, tous horizons confondus). Les valeurs en gras indiquent
  $|r| > 0{,}7$.}
\label{tab:correlations}
\begin{tabular}{lccccc}
\toprule
         & VBS   & IP     & WL     & WP    & EG     \\
\midrule
VBS      & 1,000 & 0,328  & 0,236  &-0,028 & 0,350  \\
IP       & 0,328 & 1,000  & \textbf{0,802}  & 0,541 & \textbf{0,742}  \\
WL       & 0,236 & \textbf{0,802}  & 1,000  & 0,612 & \textbf{0,741}  \\
WP       &-0,028 & 0,541  & 0,612  & 1,000 & 0,180  \\
EG       & 0,350 & \textbf{0,742}  & \textbf{0,741}  & 0,180 & 1,000  \\
\bottomrule
\end{tabular}
\end{table}

\begin{figure}[!ht]
  \centering
  \includegraphics[width=\linewidth]{figures/fig04_correlation_matrix.pdf}
  \caption{Heatmap de la matrice de corrélation inter-paramètres. La trilogie
    (IP, WL, EG) présente des corrélations $> 0{,}74$, justifiant la
    modélisation conjointe par MTGP/ICM.}
  \label{fig:correlations}
\end{figure}

% ══════════════════════════════════════════════════════════════════════════════
\section{Architecture du système et pipeline de calcul}
\label{sec:pipeline}
% ══════════════════════════════════════════════════════════════════════════════

\subsection{Infrastructure de données}

L'ensemble du pipeline s'appuie sur une base de données PostgreSQL/PostGIS
(version 17, extension PostGIS 3.4). La grille nationale est stockée dans
la table \texttt{atlas.mailles} ($N_g = 29\,407$ polygones), les prédictions
dans \texttt{atlas.ai\_interpolation\_values} ($\approx 6{,}4 \times 10^6$
valeurs pour tous modèles, paramètres et horizons). Un système de traçabilité
complet \texttt{atlas.ai\_interpolation\_runs} (310 entrées) associe chaque
ensemble de prédictions à ses méta-données : paramètres du variogramme,
métriques LOO, horodatage, version du modèle.

\subsection{Pipeline hiérarchique}

La figure~\ref{fig:pipeline} décrit l'architecture de calcul. Le pipeline
est orchestré par un worker Python (psycopg2) déclenché automatiquement
par un trigger PostgreSQL (\texttt{pg\_notify}) lors de chaque import de
nouveaux sondages, garantissant l'auto-amélioration du système.

\begin{figure}[!ht]
  \centering
  \includegraphics[width=\linewidth]{figures/fig08_pipeline_schema.pdf}
  \caption{Architecture hiérarchique du pipeline d'interpolation. Les modèles
    L1 (KED-H) et L2a (RK-SCORPAN) alimentent la Fusion BLUP (L2b). Les
    modèles L3 (VfS-PLS) et L4 (MTGP/ICM) sont complémentaires et
    indépendants. La sortie couvre 29\,407 mailles pour 5 paramètres
    × 3 horizons = 441\,105 valeurs.}
  \label{fig:pipeline}
\end{figure}

\subsection{Validation croisée Leave-One-Out}

La performance de chaque modèle est évaluée par LOO-CV, retenant successivement
un sondage en test et entraînant sur les $n-1$ restants. Le LOO-RMSE est :

\begin{equation}
  \mathrm{RMSE}_{\mathrm{LOO}} = \sqrt{\frac{1}{n}\sum_{i=1}^{n}
    \bigl[z(\bs_i) - Z^*_{(-i)}(\bs_i)\bigr]^2}
  \label{eq:loo_rmse}
\end{equation}

Pour le RK-SCORPAN, le LOO-CV inclut à la fois la ré-estimation de la
régression Ridge et le krigeage des résidus, évitant toute fuite
d'information. Pour le KED-H, les moyennes de dérive sont recalculées
à chaque itération sans le sondage retiré.

% ══════════════════════════════════════════════════════════════════════════════
\section{Résultats et validation}
\label{sec:resultats}
% ══════════════════════════════════════════════════════════════════════════════

\subsection{Performance LOO-CV comparative}
\label{subsec:loo_resultats}

Le tableau~\ref{tab:loo_rmse} présente les LOO-RMSE pour les cinq
paramètres à l'horizon H1. Sur les quinze comparaisons directes KED vs RK,
KED-H domine sur 9 critères (notamment IP, WL, WP), RK-SCORPAN sur 6
(VBS, EG et trois horizons secondaires). Cette complémentarité valide
a priori la fusion.

\begin{table}[!ht]
\centering\small
\caption{LOO-RMSE par paramètre et modèle, horizon H1.
  \textbf{Gras} = meilleur modèle par paramètre.
  La Fusion BLUP n'est pas directement évaluable par LOO-CV
  (ses poids dépendent des variances KED et RK).}
\label{tab:loo_rmse}
\begin{tabular}{lcccc}
\toprule
Paramètre  & Unité     & KED-H           & RK-SCORPAN      & VfS-PLS \\
\midrule
VBS        & g/100g    & 2,941           & \textbf{2,480}  & 2,788   \\
IP         & \%        & \textbf{9,786}  & 12,793          & —       \\
WL         & \%        & \textbf{13,213} & 15,417          & —       \\
WP         & \%        & \textbf{7,949}  & 8,071           & —       \\
EG         & \%        & 1,668           & \textbf{1,243}  & —       \\
\midrule
\multicolumn{2}{l}{Score (nb. param. gagnés)} & 9/15 & 6/15 & — \\
\bottomrule
\end{tabular}
\end{table}

La figure~\ref{fig:loo_rmse} visualise ces performances. La figure~\ref{fig:loo_scatter}
présente les diagrammes prédit vs observé pour le VBS H1, paramètre le plus
sensible à la variabilité géologique locale.

\begin{figure}[!ht]
  \centering
  \includegraphics[width=\linewidth]{figures/fig03_loo_rmse_comparison.pdf}
  \caption{\textit{Gauche :} LOO-RMSE par paramètre et modèle (H1). \textit{Droite :}
    réduction relative du RMSE RK par rapport à KED (\% positif = RK meilleur).
    VBS et EG bénéficient davantage de la covariable terrain ; IP, WL, WP
    sont mieux capturés par la dérive pédologique du KED.}
  \label{fig:loo_rmse}
\end{figure}

\begin{figure}[!ht]
  \centering
  \includegraphics[width=\linewidth]{figures/fig06_loo_scatter_vbs.pdf}
  \caption{Diagramme prédit vs observé (LOO-CV) pour le VBS à H1.
    \textit{Gauche :} KED-H ($\mathrm{RMSE}_{LOO} = 2{,}941$~g/100g).
    \textit{Droite :} RK-SCORPAN ($\mathrm{RMSE}_{LOO} = 2{,}480$~g/100g).}
  \label{fig:loo_scatter}
\end{figure}

\subsection{Dégradation de la précision avec la profondeur}

La figure~\ref{fig:multihorizon} montre l'augmentation du LOO-RMSE avec
la profondeur pour les deux modèles principaux. L'accroissement relatif
entre H1 et H3 est de 14,9\% pour KED-H et 19,0\% pour RK-SCORPAN,
reflétant la décroissance de la corrélation verticale ($r_{\mathrm{H1,H3}} = 0{,}511$).
Ce résultat est cohérent avec la nature résiduelle des sols togolais :
la minéralogie argileuse de surface (accessible aux covariables topographiques
et spectrales) est partiellement représentative de la profondeur, mais avec
une perte d'information croissante.

\begin{figure}[!ht]
  \centering
  \includegraphics[width=\linewidth]{figures/fig10_loo_multihorizon.pdf}
  \caption{LOO-RMSE du VBS par horizon de profondeur H1/H2/H3.
    La dégradation relative (+14,9\% KED, +19,0\% RK entre H1 et H3) est
    cohérente avec la corrélation verticale $r_{\mathrm{H1,H3}} = 0{,}511$
    mesurée in situ.}
  \label{fig:multihorizon}
\end{figure}

\subsection{Réduction de variance par Fusion Bayésienne}

Le tableau~\ref{tab:variance} quantifie la réduction de la variance de
krigeage induite par la Fusion BLUP. La variance combinée~\eqref{eq:blup_var}
est systématiquement inférieure aux deux modèles sources, avec une réduction
moyenne de 46,7\%~±~1,8\% par rapport à KED-H pour le VBS.

\begin{table}[!ht]
\centering\small
\caption{Variance moyenne de krigeage $\bar{\sigma}^2$ par modèle (H1).
  La réduction Fusion/KED représente le gain en certitude de la Fusion BLUP.}
\label{tab:variance}
\begin{tabular}{lcccc}
\toprule
Param. & KED-H    & RK-SCORPAN & Fusion BLUP & Réduc. \\
\midrule
VBS    & 12,45    & 10,60      & \textbf{5,73}  & 54,0\% \\
IP     & 91,0     & 79,9       & \textbf{42,3}  & 53,5\% \\
WL     & 174,2    & 140,4      & \textbf{78,2}  & 55,1\% \\
WP     & 54,5     & 60,4       & \textbf{28,1}  & 48,4\% \\
EG     & 2,92     & 2,55       & \textbf{1,31}  & 55,1\% \\
\bottomrule
\multicolumn{4}{l}{\small Unités : (g/100g)² pour VBS, (\%)² pour IP,WL,WP,EG.}
\end{tabular}
\end{table}

\begin{figure}[!ht]
  \centering
  \includegraphics[width=\linewidth]{figures/fig05_variance_reduction.pdf}
  \caption{\textit{Gauche :} variance de krigeage absolue par modèle et paramètre
    (échelle log). \textit{Droite :} réduction relative de la variance par la
    Fusion BLUP. La réduction systématique confirme la propriété théorique
    $\sigma^2_{\BLUP} < \min(\sigma^2_{\KED}, \sigma^2_{\RK})$.}
  \label{fig:variance}
\end{figure}

\subsection{Prédictions spatiales et cartographie nationale}

Le tableau~\ref{tab:predictions_grille} présente les statistiques des champs
interpolés sur la grille nationale. La figure~\ref{fig:maps} illustre la
diversité spatiale des prédictions VBS selon les quatre modèles.

\begin{table}[!ht]
\centering\small
\caption{Statistiques des prédictions sur la grille nationale (H1, 29\,407 mailles).}
\label{tab:predictions_grille}
\begin{tabular}{llccc}
\toprule
Param. & Modèle     & $\bar{Z}^*$ & $\sigma_{Z^*}$ & Unité \\
\midrule
VBS & KED-H       & 3,65  & 1,36  & g/100g \\
VBS & RK-SCORPAN  & 3,80  & 2,20  & g/100g \\
VBS & Fusion BLUP & 3,83  & 1,48  & g/100g \\
IP  & KED-H       & 18,99 & 4,58  & \%     \\
IP  & RK-SCORPAN  & 19,00 & 6,61  & \%     \\
WL  & KED-H       & 36,07 & 5,00  & \%     \\
WL  & RK-SCORPAN  & 38,10 & 11,76 & \%     \\
WP  & KED-H       & 17,45 & 3,05  & \%     \\
EG  & KED-H       & 3,82  & 0,82  & \%     \\
EG  & RK-SCORPAN  & 2,99  & 1,93  & \%     \\
\bottomrule
\end{tabular}
\end{table}

\begin{figure}[!ht]
  \centering
  \includegraphics[width=\linewidth]{figures/fig07_maps_vbs_h1.pdf}
  \caption{Cartes de prédiction du VBS à H1 selon les modèles L1 à L3.
    KED-H montre une interpolation lisse dominée par la structure pédologique ;
    RK-SCORPAN révèle des gradients topographiques N-S ; la Fusion combine
    les deux signaux avec incertitude réduite.}
  \label{fig:maps}
\end{figure}

\subsection{Résultats VfS-PLS (Sentinel-2)}

La régression PLS à 3 composantes sur les indices spectraux
$\{I_{\mathrm{argile}}, I_{\mathrm{SWIR}}, \mathrm{NDVI}, I_{\mathrm{Fe}}\}$
produit un LOO-RMSE de 2,788~g/100g pour le VBS H1, inférieur au
KED-H (2,941~g/100g). Cette performance est remarquable compte tenu
du paradoxe de profondeur : Sentinel-2 ne pénètre que les 2~mm superficiels,
alors que les mesures de VBS concernent le sol à 0–1~m. Sa performance
s'explique par la corrélation verticale partielle ($r_{\mathrm{H1,H3}} = 0{,}511$)
dans les sols résiduels togolais (53\% du territoire).

Sur les 29\,407 mailles nationales, 24\,038 (81,7\%) reçoivent une prédiction
VfS valide après masquage des cuirasses latéritiques ($N = 2\,464$ mailles,
8,4\%) et des alluvions récentes (profil vertical discontinu).
La moyenne nationale prédite ($\bar{z}_{\mathrm{VfS}} = 3{,}53$~g/100g)
est cohérente avec la mesure in situ ($\bar{z}_{\mathrm{obs}} = 4{,}27$~g/100g,
biais attendu de la couche superficielle vs. profondeur 0–1~m).

\subsection{Profil de prédiction et incertitude spatiale}

La figure~\ref{fig:transect} présente le transect méridien (N-S) des
prédictions VBS avec leurs intervalles de confiance à 95\%. Les trois
modèles présentent des profils divergents dans les zones de transition
géologique, notamment entre la région centrale (gneiss, VBS modéré)
et la région de Kara (schistes, VBS plus variable). La Fusion BLUP
produit l'enveloppe d'incertitude la plus étroite dans toutes les zones,
confirmant la réduction de variance quantifiée en section~\ref{subsec:loo_resultats}.

\begin{figure}[!ht]
  \centering
  \includegraphics[width=\linewidth]{figures/fig09_transect_uncertainty.pdf}
  \caption{Transect méridien N-S (6°N--11°N) des prédictions VBS H1.
    Les zones colorées représentent les intervalles de confiance à 95\%
    ($\pm 2\sigma$). Les régions physiographiques sont indiquées en abscisse.
    La Fusion BLUP (vert) maintient l'enveloppe d'incertitude la plus
    étroite dans toutes les zones.}
  \label{fig:transect}
\end{figure}

% ══════════════════════════════════════════════════════════════════════════════
\section{Discussion}
\label{sec:discussion}
% ══════════════════════════════════════════════════════════════════════════════

\subsection{Interprétation géotechnique des résultats}

La distribution spatiale du VBS prédit (tableau~\ref{tab:predictions_grille})
révèle une hétérogénéité géotechnique nationale significative, cohérente
avec la minéralogie attendue. Les zones de VBS élevé ($> 6$~g/100g) se
concentrent dans les alluvions de la région maritime (smectites héritées
de la sédimentation lacustre) et dans les Vertisols de la Dépression de
la Lama (argiles 2:1 gonflantes) \cite{chassagneux1996}.

Les gradients N-S capturés par RK-SCORPAN, absents du champ KED-H, correspondent
à des transitions climatiques réelles : les précipitations annuelles décroissent
de 1\,400~mm (côte atlantique) à 800~mm (nord sahélien), modifiant le degré
de latérisation et donc l'activité argileuse \cite{brgm2019togo}. Cette
composante longue portée est la contribution spécifique du modèle SCORPAN.

\subsection{Position dans la littérature}

Trois innovations méthodologiques distinguent ce travail :

\begin{enumerate}
  \item \textbf{Dérive géologique hiérarchique à cinq niveaux} : La combinaison
    zones géomorphologiques + pédologie + risque RGA + géologie + hydrogéologie
    comme variable de dérive structurée n'est pas documentée en Afrique de l'Ouest.
    Elle améliore significativement le KED classique \cite{odeh1994further}
    dans les zones à faible densité de sondages.

  \item \textbf{VBS par télédétection (VfS)} : À notre connaissance, aucune
    étude publiée ne relie les indices spectraux Sentinel-2 à la Valeur au
    Bleu de Méthylène pour des sols tropicaux togolais \cite{castaldi2019sentinel}.
    La performance LOO-RMSE = 2,788~g/100g surpassant KED-H (2,941~g/100g)
    établit une preuve de concept pour la géotechnique sans laboratoire.

  \item \textbf{Fusion BLUP adaptative} : La combinaison de KED et RK par
    pondération inverse à l'incertitude locale est une adaptation originale
    de la théorie BLUP \cite{chilesdelfiner2012} à la géotechnique tropicale.
    La réduction de variance systématique (46,7\%~±~1,8\% sur le VBS)
    dépasse les résultats de la fusion naïve (50/50) qui produirait une
    réduction théorique de 38\% seulement.
\end{enumerate}

\subsection{Limites et perspectives}

Trois limites méritent d'être soulignées. Premièrement, la densité de sondages
($\rho = 2{,}2 \times 10^{-3}$~km$^{-2}$) reste insuffisante pour caractériser
les hétérogénéités à courte portée dans les zones géologiques spéciales
(Dépression de la Lama, Fosse aux Lions, Plaine du Mono). L'auto-amélioration
du système sur import de nouveaux sondages est activée par un trigger
PostgreSQL dès $n \geq 10$ sondages par paramètre, mais la précision
du variogramme restera limitée tant que $n < 150$.

Deuxièmement, le modèle CatBoost (L3 ML pur), prévu dans la hiérarchie,
n'est pas activé car $n < 200$, seuil minimal pour éviter le sur-apprentissage
avec validation spatiale en blocs. Son déploiement est conditionné à
l'importation de 80 sondages supplémentaires, estimés disponibles dans
les archives de la Direction des Infrastructures du Togo.

Troisièmement, la performance VfS-PLS est limitée aux sols résiduels
(53\% du territoire). Pour les zones alluvionnaires et les cuirasses
latéritiques (19\% du territoire), la prédiction spectrale est signalée
avec un indicateur d'incertitude haute.

\subsection{Gain du co-krigeage MTGP}

La figure~\ref{fig:mtgp} détaille le gain apporté par le MTGP/ICM
par rapport au krigeage indépendant par paramètre. Le gain le plus
important concerne EG ($-$12,7\%), dont l'échantillonnage limité
($n = 64$) est compensé par les corrélations avec IP ($r = 0{,}742$)
et WL ($r = 0{,}741$), tous deux mieux documentés ($n \geq 93$).

\begin{figure}[!ht]
  \centering
  \includegraphics[width=\linewidth]{figures/fig13_mtgp_gain.pdf}
  \caption{\textit{Gauche :} LOO-RMSE mono-krigeage vs MTGP/ICM par
    paramètre (H1). \textit{Droite :} réduction relative RMSE par le
    co-krigeage. EG bénéficie le plus grâce à ses fortes corrélations
    avec IP et WL.}
  \label{fig:mtgp}
\end{figure}

\subsection{Synthèse de la hiérarchie L1--L4}

La figure~\ref{fig:synthesis} présente une synthèse multi-critères
de la hiérarchie. Le graphe radar compare les performances relatives
normalisées; le panel B quantifie la couverture spatiale de chaque
modèle; le panel C visualise la réduction de variance par Fusion BLUP.

\begin{figure}[!ht]
  \centering
  \includegraphics[width=\linewidth]{figures/fig14_synthesis.pdf}
  \caption{Synthèse comparative — hiérarchie L1--L4.
    \textit{A :} radar performances normalisées.
    \textit{B :} couverture spatiale (×1000 mailles).
    \textit{C :} réduction variance Fusion BLUP.}
  \label{fig:synthesis}
\end{figure}

% ══════════════════════════════════════════════════════════════════════════════
\section{Conclusion}
\label{sec:conclusion}
% ══════════════════════════════════════════════════════════════════════════════

La hiérarchie de méthodes développée constitue, à la connaissance de l'auteur,
la première application rigoureusement validée de la Régression Kriging SCORPAN
à la prédiction de paramètres géotechniques en Afrique subsaharienne.
L'originalité technique réside dans la combinaison de cinq innovations
complémentaires --- dérive hiérarchique, validation LOO complète du RK,
Fusion Bayésienne BLUP, prédiction spectrale VfS, co-krigeage MTGP ---
chacune apportant une contribution incrémentale mesurable.

Cet article présente la première cartographie géotechnique nationale du Togo
à haute résolution spatiale (2~km~×~2~km), fondée sur une hiérarchie de modèles
géostatistiques rigoureusement validés. Les contributions principales sont :

\begin{enumerate}
  \item Un Krigeage avec Dérive Hiérarchique à cinq niveaux géologiques
    (KED-H), produisant $441\,105$ valeurs interpolées avec variance quantifiée ;
  \item Une Régression Kriging SCORPAN complète (RK-SCORPAN) avec validation
    LOO rigoureuse (ré-estimation régression + krigeage à chaque itération) ;
  \item Une Fusion Bayésienne BLUP réduisant la variance de 46,7\%~±~1,8\%
    par rapport au modèle unilatéral le plus précis ;
  \item Une preuve de concept de prédiction géotechnique par télédétection
    satellitaire (VfS-PLS, Sentinel-2) avec LOO-RMSE = 2,788~g/100g
    pour le VBS, inférieur au krigeage géostatistique classique ;
  \item Un Processus Gaussien Multi-Tâches (MTGP/ICM) exploitant les
    corrélations entre les cinq paramètres géotechniques pour améliorer
    la précision sur les paramètres peu documentés (EG, WP).
\end{enumerate}

L'atlas géotechnique numérique résultant, stocké dans une base de données
PostgreSQL/PostGIS et accessible via une API REST, constitue un outil
opérationnel pour le dimensionnement des infrastructures civiles au Togo.
Son architecture auto-améliorante garantit que chaque import de nouveaux
sondages déclenche automatiquement le recalcul des modèles et la mise à
jour des cartes, assurant la pérennité et la pertinence du système à long terme.

% ══════════════════════════════════════════════════════════════════════════════
\section{Analyse de sensibilité et robustesse}
\label{sec:sensibilite}
% ══════════════════════════════════════════════════════════════════════════════

\subsection{Sensibilité au modèle variographique}

Le choix du modèle variographique (sphérique, exponentiel, gaussien) constitue
une source d'incertitude épistémique dans le processus de krigeage. Une analyse
de sensibilité est conduite pour le VBS H1 en comparant trois modèles :

\begin{table}[!ht]
\centering\small
\caption{Sensibilité du LOO-RMSE KED-H au modèle variographique (VBS H1).
  Les paramètres sont recalibrés par moindres carrés pondérés pour chaque modèle.}
\label{tab:sensibilite_variogram}
\begin{tabular}{lcccc}
\toprule
Modèle $\gamma(h)$ & $C_0$ & $C_0{+}C$ & $a$ (km) & RMSE \\
\midrule
Sphérique (retenu)   & 2,80 & 14,0 & 220 & 2,941 \\
Exponentiel          & 2,65 & 13,5 & 185 & 2,958 \\
Gaussien             & 3,10 & 14,8 & 175 & 3,012 \\
Matérn 3/2           & 2,72 & 13,8 & 198 & 2,947 \\
\bottomrule
\end{tabular}
\end{table}

La différence de LOO-RMSE entre les modèles est inférieure à 2,5\% (0,07~g/100g),
indiquant une robustesse élevée aux choix du modèle variographique. Cette
robustesse est caractéristique des données géotechniques à portée longue ($a > 150$~km)
où la structure du variogramme à courte portée est peu contrainte par les
données de validation croisée \cite{lark2000kriging}.

\subsection{Sensibilité à la taille de l'échantillon}

La figure~\ref{fig:sensib_n} illustre l'évolution du LOO-RMSE KED-H en fonction
de la taille de l'échantillon $n$ pour le VBS H1, obtenue par sous-échantillonnage
aléatoire répété ($B = 50$ répétitions). L'analyse révèle une courbe d'apprentissage
typique : le LOO-RMSE décroît rapidement jusqu'à $n \approx 60$, puis se stabilise
avec un gain marginal décroissant.

La relation empirique ajustée est :
\begin{equation}
  \mathrm{RMSE}_{\mathrm{LOO}}(n) \approx a_0 \exp\!\left(-\frac{n}{n_0}\right) + \mathrm{RMSE}_\infty
  \label{eq:learning_curve}
\end{equation}
avec $a_0 = 4{,}8$~g/100g, $n_0 = 28$ (constante de saturation),
$\mathrm{RMSE}_\infty = 2{,}75$~g/100g (plancher théorique lié à la pépite
et à la variabilité intrinsèque du sol).

\begin{figure}[!ht]
  \centering
  \includegraphics[width=\linewidth]{figures/fig15_learning_curve.pdf}
  \caption{Courbe d'apprentissage LOO-RMSE pour KED-H et RK-SCORPAN
    en fonction de $n$ (sous-échantillonnage aléatoire, $B = 50$
    répétitions). La ligne pointillée verticale indique $n = 95$
    (échantillon actuel). Le plancher $\mathrm{RMSE}_\infty$ est la
    borne inférieure théorique liée à l'effet de pépite.}
  \label{fig:learning_curve}
\end{figure}

Cette relation prédit que doubler l'échantillon ($n = 95 \to 190$) réduirait
le LOO-RMSE de $\sim$0,12~g/100g, une amélioration modeste confirmant que
l'architecture actuelle est proche de la saturation de performance pour la
résolution de 2~km. Cette saturation est cohérente avec la portée du variogramme
($a = 220$~km) : à la résolution de 2~km, le krigeage ne peut pas résoudre
les hétérogénéités à échelle inférieure à $a/n^{1/2} \approx 22$~km,
indépendamment de la densité de sondages.

\subsection{Influence du paramètre de régularisation sur la prédiction spatiale}

Pour quantifier l'effet de $\lambda_R$ sur les cartes prédites (et non
seulement sur le LOO-RMSE de la régression seule), une analyse en grille
complète ($N_g = 29\,407$ mailles) est conduite avec trois valeurs de
$\lambda_R \in \{0{,}05, 0{,}18, 1{,}0\}$ :

\begin{table}[!ht]
\centering\small
\caption{Sensibilité des prédictions RK-SCORPAN sur la grille nationale
  à $\lambda_R$ (VBS H1). $\bar{Z}^*$ : moyenne prédite sur la grille,
  $\sigma_{Z^*}$ : écart-type spatial, $\Delta\sigma^2$ : variation de
  la variance de prédiction.}
\label{tab:sensibilite_ridge}
\begin{tabular}{ccccc}
\toprule
$\lambda_R$ & $\bar{Z}^*$ & $\sigma_{Z^*}$ & RMSE$_{\mathrm{LOO}}$ & $\Delta\sigma^2$ \\
\midrule
0,05  & 3,82 & 2,41 & 2,531 & $+$8,3\% \\
0,18* & 3,80 & 2,20 & 2,480 & 0 (réf.) \\
1,00  & 3,74 & 1,85 & 2,612 & $-$5,8\% \\
\bottomrule
\multicolumn{5}{l}{\small * Valeur retenue. Unités : g/100g.}
\end{tabular}
\end{table}

La valeur $\lambda_R = 0{,}05$ (sous-régularisation) amplifie les
gradients spatiaux ($\sigma_{Z^*} = 2{,}41$ vs $2{,}20$ à l'optimum),
produisant des artefacts de prédiction dans les zones extrapolées
(hors de la convexe hull des sondages). À l'inverse, $\lambda_R = 1{,}0$
(sur-régularisation) écrase les gradients, réduisant l'écart-type spatial
de 16\% et sous-estimant les zones extrêmes (VBS $>$ 8~g/100g).

\subsection{Performance par région administrative}

Le tableau~\ref{tab:perf_region} présente les LOO-RMSE KED-H et RK-SCORPAN
décomposés par région administrative (niveau ADM1), permettant d'identifier
les zones géographiques où l'interpolation est la plus incertaine.

\begin{table}[!ht]
\centering\small
\caption{LOO-RMSE VBS H1 par région administrative (ADM1) et modèle.
  $n$ : nombre de sondages dans la région. Les régions avec $n < 10$
  sont exclues de la statistique LOO régionale.}
\label{tab:perf_region}
\begin{tabular}{lccc}
\toprule
Région & $n$ & KED-H & RK-SCORPAN \\
\midrule
Maritime    & 38 & 2,410 & 2,182 \\
Plateaux    & 22 & 3,218 & 2,891 \\
Centrale    & 18 & 2,847 & 2,531 \\
Kara        & 14 & 3,412 & 3,105 \\
Savanes     & 10 & 3,685 & 3,248 \\
\midrule
National    & 102& 2,941 & 2,480 \\
\bottomrule
\end{tabular}
\end{table}

Les régions septentrionales (Kara, Savanes) présentent les erreurs les plus
élevées, cohérentes avec la faible densité de sondages et la géologie complexe
(alternance schistes/quartzites/grès voltaïens). La région Maritime, bénéficiant
du plus grand nombre de sondages ($n = 38$) et d'une géologie plus homogène
(alluvions quaternaires), présente les meilleures performances pour les deux modèles.

\begin{figure}[!ht]
  \centering
  \includegraphics[width=\linewidth]{figures/fig16_regional_performance.pdf}
  \caption{\textit{Gauche :} LOO-RMSE VBS H1 par région administrative.
    Les traits pointillés horizontaux indiquent les moyennes nationales.
    \textit{Droite :} variance de krigeage KED-H moyenne sur la grille
    par région (indicateur d'incertitude spatiale).}
  \label{fig:regional}
\end{figure}

% ══════════════════════════════════════════════════════════════════════════════
\section{Calcul des covariables SCORPAN}
\label{sec:covariables}
% ══════════════════════════════════════════════════════════════════════════════

\subsection{Modèle numérique de terrain (DSM)}

La topographie est la covariable terrain la plus informative du modèle SCORPAN.
À partir du Modèle Numérique de Surface (DSM) SRTM 30~m, quatre dérivés sont
calculés sur la grille nationale :

\paragraph{Altitude moyenne} ($z_{\mathrm{DSM}}$, m) : moyenne des valeurs
DSM dans la maille de 2~km~$\times$~2~km. Contrôle le degré de latérisation
et la perméabilité du substrat.

\paragraph{Pente} ($\alpha$, \%): calculée par différences finies centrées
sur une fenêtre 3~$\times$~3 :
\begin{equation}
  \alpha = \arctan\!\left(\sqrt{\left(\frac{\partial z}{\partial x}\right)^2
    + \left(\frac{\partial z}{\partial y}\right)^2}\right) \times \frac{180}{\pi}
  \label{eq:slope}
\end{equation}
Une pente élevée corrèle négativement avec la teneur en argile fine
(lessivage mécanique) et donc avec VBS et IP.

\paragraph{Indice de Position Topographique (TPI)} : différence entre
l'altitude d'un point et la moyenne de son voisinage circulaire de rayon
$R = 5$~km :
\begin{equation}
  \mathrm{TPI}(\bs) = z(\bs) - \frac{1}{|\mathcal{N}_R|}\sum_{\bs' \in \mathcal{N}_R} z(\bs')
  \label{eq:tpi}
\end{equation}
Les TPI négatifs (fonds de vallée, plaines) indiquent des zones d'accumulation
argileuse (VBS élevé) ; les TPI positifs (interfluves) des zones de drainage.

\paragraph{HAND (Height Above Nearest Drainage)} : hauteur verticale
au-dessus du réseau hydrographique le plus proche, calculée par algorithme
D8 de routage de flux :
\begin{equation}
  \mathrm{HAND}(\bs) = z(\bs) - z_{\mathrm{rivière}}(\bs)
  \label{eq:hand}
\end{equation}
Cette variable encode le potentiel de saturation hydrique, facteur
critique pour EG (Coefficient de Gonflement) \cite{goovaerts1997}.

\subsection{Covariables climatiques (WorldClim)}

Les données WorldClim v2.1 à résolution 30'' (environ 1~km) fournissent :

\begin{itemize}
  \item $P_{\mathrm{ann}}$ (mm) : précipitation annuelle totale (BIO12).
    Contrôle le degré d'altération et la profondeur des argiles.
  \item $P_{\mathrm{sèche}}$ (mm) : précipitation du trimestre le plus
    sec, calculée comme $P_{\mathrm{sèche}} = \sum_{\text{trim. sec}} P_i$.
    Corrèle avec le retrait-gonflement saisonnier (EG).
  \item $\sigma_P$ (mm) : saisonnalité pluviométrique (BIO15, coefficient
    de variation mensuel). Les sols soumis à une forte saisonnalité
    développent des fentes de retrait (IP et WL élevés).
\end{itemize}

Ces trois covariables sont calculées par intersection spatiale PostGIS
entre les rasters WorldClim et la grille nationale, avec agrégation
par moyenne pondérée sur la surface de chaque maille.

\subsection{Covariables géologiques et pédologiques}

Les variables catégorielles (type de sol, formation géologique, classe
de risque RGA) sont encodées en variables ordinales selon leur activité
argileuse croissante :

\begin{equation}
  \mathrm{enc}(c) = \frac{\bar{z}_c - \bar{z}_{\mathrm{nat}}}{\sigma_{\mathrm{nat}}}
  \label{eq:encoding}
\end{equation}

où $\bar{z}_c$ est la moyenne du paramètre dans la catégorie $c$, et
$\bar{z}_{\mathrm{nat}}$, $\sigma_{\mathrm{nat}}$ sont la moyenne et
l'écart-type nationaux. Cet encodage numérique préserve le sens physique
(contraste entre formations géologiques) et améliore la convergence
de la régression Ridge \eqref{eq:ridge}.

\subsection{Résumé des covariables}

Le tableau~\ref{tab:covariables} liste les 12 covariables utilisées dans
le modèle RK-SCORPAN, leur source et leur pertinence géotechnique.

\begin{table}[!ht]
\centering\small
\caption{Covariables du modèle RK-SCORPAN. Catégorie SCORPAN :
  $S$~=~sol, $C$~=~climat, $R$~=~relief, $P$~=~matériau parental,
  $N$~=~position.}
\label{tab:covariables}
\setlength{\tabcolsep}{4pt}
\begin{tabular}{llcp{2.8cm}}
\toprule
Variable & Source & Cat. & Pertinence géotech. \\
\midrule
Altitude DSM       & SRTM 30m     & R & Latérisation \\
Pente ($\alpha$)   & SRTM dérivé  & R & Lessivage mécanique \\
TPI (5~km)        & SRTM dérivé  & R & Accumulation argileuse \\
HAND               & SRTM dérivé  & R & Saturation hydrique \\
$P_{\mathrm{ann}}$ & WorldClim    & C & Altération chimique \\
$P_{\mathrm{sèche}}$& WorldClim   & C & Retrait-gonflement \\
$\sigma_P$         & WorldClim    & C & Plasticité saisonnière \\
Géologie (enc.)   & BRGM 2019    & P & Minéralogie héritée \\
Pédologie (enc.)  & FAO          & S & Type d'argile \\
Risque RGA (enc.) & Chassagneux  & S & Activité gonflante \\
Coord. Est (UTM)  & --           & N & Gradient W-E \\
Coord. Nord (UTM) & --           & N & Gradient S-N \\
\bottomrule
\end{tabular}
\end{table}

% ══════════════════════════════════════════════════════════════════════════════
\section{Implémentation algorithmique}
\label{sec:algo}
% ══════════════════════════════════════════════════════════════════════════════

\subsection{Algorithme KED-H}

L'algorithme KED-H est implémenté en Python avec la bibliothèque PyKrige
(v1.7) pour la factorisation de la matrice de covariance. La complexité
algorithmique est dominée par la résolution du système~\eqref{eq:ked_system} :
$\mathcal{O}(n^3)$ pour la factorisation, $\mathcal{O}(n^2 N_g)$ pour
la prédiction sur $N_g$ mailles.

\begin{description}
  \item[\textbf{Étape 1 --- Prior hiérarchique.}] Calculer $\mu_h(\bs_i)$ selon
    \eqref{eq:prior_hierarchique} pour tous les sondages et toutes les mailles
    de la grille. En cas de catégorie sans mesure ($n_c < 5$), repli
    automatique au niveau supérieur de la hiérarchie.

  \item[\textbf{Étape 2 --- Résidus de dérive.}] Calculer les résidus
    $\varepsilon_i = z_i - \mu_h(\bs_i)$ aux sondages. Ces résidus
    sont présumés stationnaires de second ordre et alimentent le
    calcul du variogramme empirique.

  \item[\textbf{Étape 3 --- Ajustement variographique.}] Calculer le
    semi-variogramme empirique omnidirectionnel par la méthode des moments :
    \begin{equation}
      \hat{\gamma}(h \pm \Delta h) = \frac{1}{2|N(h)|}\sum_{(i,j)\in N(h)}
        [\varepsilon(\bs_i) - \varepsilon(\bs_j)]^2
      \label{eq:gamma_emp}
    \end{equation}
    où $N(h) = \{(i,j) : |\bs_i - \bs_j| \in [h - \Delta h, h + \Delta h]\}$.
    Ajustement du modèle sphérique \eqref{eq:variogram_spherique} par
    moindres carrés pondérés ($\propto |N(h)|$).

  \item[\textbf{Étape 4 --- Krigeage.}] Résoudre \eqref{eq:ked_system}
    pour chaque maille $\bs_0$. La matrice $\bK$ est assemblée une seule
    fois (taille $n \times n$) et factorisée par Cholesky
    ($\mathcal{O}(n^3)$). Les prédictions sont effectuées par substitution
    avant-arrière ($\mathcal{O}(n^2)$ par maille).

  \item[\textbf{Étape 5 --- Traçabilité.}] Enregistrer dans
    \texttt{ai\_interpolation\_runs} : run\_id, paramètres du variogramme,
    $n_{\mathrm{train}}$, LOO-RMSE, durée de calcul, version du modèle.
\end{description}

\subsection{Validation LOO-CV rigoureuse pour le RK}

La validation croisée LOO du modèle RK exige une précaution particulière :
à chaque itération $i$, la régression Ridge \emph{et} le krigeage des
résidus doivent être ré-estimés sans le sondage $\bs_i$. L'omission
de cette ré-estimation conduit à un optimisme de l'ordre de 15--30\%
sur le LOO-RMSE. L'algorithme est :

\begin{description}
  \item[\textbf{Pour} $i = 1, \ldots, n$ :]
    \begin{enumerate}
      \item Estimer $\hat{\bm{\beta}}^{(-i)}$ par Ridge sur
        $\{\bs_j, z_j\}_{j \neq i}$ ;
      \item Calculer les résidus $\hat{r}_j^{(-i)} = z_j - \mathbf{x}_j^T
        \hat{\bm{\beta}}^{(-i)}$ pour $j \neq i$ ;
      \item Ajuster le variogramme des résidus $\hat{\gamma}^{(-i)}(h)$
        sur $n-1$ points ;
      \item Prédire par krigeage ordinaire en $\bs_i$ :
        $\hat{r}^{(-i)}(\bs_i)$ ;
      \item Prédiction LOO :
        $Z^*_{(-i)}(\bs_i) = \mathbf{x}_i^T \hat{\bm{\beta}}^{(-i)} + \hat{r}^{(-i)}(\bs_i)$.
    \end{enumerate}
\end{description}

Cette procédure, de complexité $\mathcal{O}(n^4)$ dans le pire cas,
est rendue praticable par la taille modeste de l'échantillon ($n \leq 95$
par paramètre et horizon), avec un temps de calcul typique de 12--45~minutes
par paramètre.

\subsection{Calibration du paramètre Ridge $\lambda_R$}

Le paramètre de régularisation Ridge est calibré par LOO-CV interne
sur les sondages d'entraînement. La figure~\ref{fig:ridge_cv} illustre
la courbe de calibration pour le VBS H1, ainsi que les chemins de
régularisation des 12 coefficients.

\begin{figure}[!ht]
  \centering
  \includegraphics[width=\linewidth]{figures/fig11_ridge_calibration.pdf}
  \caption{\textit{Gauche :} courbe de calibration LOO-RMSE de la régression
    Ridge en fonction de $\lambda_R$. \textit{Droite :} chemins de régularisation
    des 12 coefficients normalisés. La ligne rouge indique $\lambda^*_R = 0{,}18$.}
  \label{fig:ridge_cv}
\end{figure}

Le critère de sélection est le LOO-RMSE de la régression seule (sans krigeage
des résidus), évalué sur une grille logarithmique $\lambda_R \in [10^{-4}, 10^2]$.
Pour le VBS H1, la valeur optimale est $\lambda_R^* = 0{,}18$, correspondant
à un compromis biais-variance adapté aux 12 covariables et $n = 95$ observations.

\subsection{Analyse des composantes PLS (VfS)}

Les trois composantes PLS retenues pour le modèle VfS expliquent respectivement
68,5\%, 18,3\% et 8,2\% de la variance des indices spectraux (Fig.~\ref{fig:pls}).
La première composante est dominée par $I_{\mathrm{argile}} = B_{11}/B_{12}$
(chargement~$= 0{,}72$), confirmant son rôle de marqueur principal de l'activité
argileuse. La deuxième composante capture l'opposition NDVI/$I_{\mathrm{SWIR}}$,
reflétant le gradient sol nu/végétation qui module la réflectance de surface.

\begin{figure}[!ht]
  \centering
  \includegraphics[width=\linewidth]{figures/fig12_pls_loadings.pdf}
  \caption{\textit{Gauche :} chargements PLS des 3 composantes latentes
    sur les 4 indices spectraux Sentinel-2. CP1 : indice d'activité argileuse ;
    CP2 : gradient sol/végétation ; CP3 : SWIR résiduel.
    \textit{Droite :} variance expliquée dans $\mathbf{X}$ (indices) et
    $y$ (VBS).}
  \label{fig:pls}
\end{figure}

\subsection{Paramètres variographiques par paramètre}

Le tableau~\ref{tab:variogrammes} rassemble les paramètres du modèle
variographique sphérique ajusté pour chaque paramètre à l'horizon H1.
L'effet de pépite relatif ($C_0/(C_0+C)$) quantifie la part de variabilité
non capturée à la résolution de la grille ; une valeur élevée indique
une variabilité à courte portée qui limite la prédiction spatiale.

\begin{table}[!ht]
\centering\small
\caption{Paramètres du modèle variographique sphérique ajusté (H1).
  $C_0$ : effet de pépite, $C$ : variance structurale, $a$ : portée.
  Rapport de pépite : $C_0/(C_0+C)$.}
\label{tab:variogrammes}
\begin{tabular}{lcccc}
\toprule
Param. & $C_0$  & $C_0+C$ & $a$ (km) & $C_0/(C_0+C)$ \\
\midrule
VBS    & 2,80   & 14,00   & 220      & 0,200 \\
IP     & 22,0   & 145,0   & 310      & 0,152 \\
WL     & 45,0   & 185,0   & 280      & 0,243 \\
WP     & 15,0   & 72,0    & 195      & 0,208 \\
EG     & 0,85   & 3,60    & 175      & 0,236 \\
\bottomrule
\multicolumn{5}{l}{\small Unités: (g/100g)² pour VBS, (\%)² pour IP,WL,WP,EG.}
\end{tabular}
\end{table}

Les portées importantes ($175$--$310$~km) reflètent la structure
géologique à grande échelle du Togo : les cinq grandes régions
physiographiques (Maritime, Plateaux, Centrale, Kara, Savanes) présentent
une continuité spatiale sur des centaines de kilomètres. Les effets de pépite
modérés (0,15--0,24) indiquent une variabilité à courte portée réelle
(hétérogénéités locales de sol) qui ne peut être résolue à la résolution
de 2~km.

% ══════════════════════════════════════════════════════════════════════════════
\section{Analyse géologique des résultats}
\label{sec:geologie}
% ══════════════════════════════════════════════════════════════════════════════

\subsection{Distribution spatiale du VBS et contrôles géologiques}

\clearpage

La carte de prédiction VBS (Fig.~\ref{fig:maps}) révèle trois régimes
géotechniques distincts, cohérents avec la géologie du substrat :

\paragraph{Régime argileux (VBS > 4,5~g/100g, 22\% du territoire).}
Concentré dans la région maritime (alluvions quaternaires lacustres à
smectites dominantes), la Dépression de la Lama (argiles 2:1 gonflantes,
VBS mesuré moyen~= 4,11~g/100g) et les plaines alluviales du Mono.
Ces formations sont caractérisées par un IP élevé ($> 30\%$) et un
potentiel de gonflement fort à très fort.

\paragraph{Régime intermédiaire (VBS 2,5--4,5~g/100g, 53\% du territoire).}
Dominant sur les gneiss et orthogneiss du socle cristallin, où l'altération
produit un mélange kaolinite-illite. C'est le régime géotechnique le plus
répandu au Togo, avec des comportements de sol prévisibles par les modèles
géostatistiques (LOO-RMSE KED-H = 2,941~g/100g).

\paragraph{Régime faiblement argileux (VBS < 2,5~g/100g, 25\% du territoire).}
Associé aux quartzites et grès de l'Atacora (nord-ouest) et aux grès
voltaïens (extrême nord). Ces formations produisent des sols sableux à
faible activité argileuse, bien prédits par le VfS-PLS sur la couche
superficielle (kaolinite prédominante, spectre SWIR caractéristique).

\subsection{Comparaison zones géologiques spéciales vs zones générales}

La Dépression de la Lama (zone Z1, 5\% du territoire) est la zone la mieux
documentée des 5 zones spéciales : 12 sondages y sont disponibles, permettant
une estimation locale fiable. Le VBS moyen mesuré y est de 4,11~g/100g
avec un écart-type de 3,61~g/100g, reflétant une forte hétérogénéité interne
cohérente avec la nature lacustre de sa sédimentation (alternance de couches
argileuses à smectite et de couches limono-sableuses).

À l'inverse, trois zones spéciales (Oti, Mono, Fosse aux Lions) ne contiennent
aucun sondage dans la base de données actuelle. Les prédictions dans ces zones
reposent entièrement sur l'interpolation à partir des sondages voisins et
sur les moyennes pédologiques/géologiques du prior hiérarchique. La variance
de krigeage y est typiquement 2 à 3 fois supérieure à la variance nationale
moyenne, ce qui se traduit par des intervalles de prédiction élargis
(IC 95\% couvrant $\pm$~4--6~g/100g de VBS).

\subsection{Anomalies géotechniques identifiées}

L'analyse comparative KED-H vs RK-SCORPAN révèle des \emph{anomalies
positives} (RK prédit VBS plus élevé que KED) dans deux zones :

\begin{enumerate}
  \item \textbf{Région de Sotouboua (centrale)} : Le modèle RK prédit
    VBS~=~5,8~g/100g vs KED~=~3,9~g/100g. Cette zone correspond à
    des shales et argiles sédimentaires du Voltaïen, dont la haute
    teneur en illites est mal représentée dans la dérive pédologique
    du KED (trop peu de sondages dans cette formation : $n = 3$).
    La Fusion BLUP attribue un poids plus élevé au RK dans cette zone
    ($\beta \approx 0{,}72$).

  \item \textbf{Couloir Togoville--Vogan (maritime)} : VBS~=~7,2~g/100g
    (RK) vs 5,1~g/100g (KED). Correspond à la zone de transition
    lagune-continent avec des argiles lacustres peu représentées dans
    l'échantillonnage actuel (2 sondages seulement). Zone prioritaire
    pour la campagne terrain future.
\end{enumerate}

\subsection{Relation VBS--IP : droite de régression et cohérence physique}

La corrélation VBS--IP ($r = 0{,}328$) est plus faible qu'attendu physiquement
(la smectite produit simultanément VBS et IP élevés). Cette sous-corrélation
reflète la confusion entre smectites (forte activité) et illites (activité
intermédiaire) dans l'échantillonnage. La droite de régression observée est :

\begin{equation}
  \mathrm{IP} = 4{,}72 + 3{,}93 \times \mathrm{VBS}, \quad
  R^2 = 0{,}107, \quad n = 310
  \label{eq:vbs_ip_regression}
\end{equation}

Cette relation, bien que faible à l'échelle nationale, devient significative
($R^2 = 0{,}38$, $p < 0{,}001$) au sein des Vertisols (argiles 2:1 dominantes),
validant l'utilisation du VBS comme proxy partiel de l'IP dans les modèles MTGP.

\subsection{Interprétation de la matrice de coregionalisation}

La matrice de coregionalisation $\hat{\mathbf{B}}$ estimée par MTGP/ICM
(annexe~\ref{app:mtgp}, équation~\eqref{eq:coregionalisation}) s'interprète
physiquement comme la matrice de corrélation entre les cinq paramètres à
portée infinie (pour $h \to 0$). Les valeurs élevées de $B_{23} = 0{,}80$
(IP, WL) et $B_{24} = 0{,}54$ (IP, WP) reflètent les contraintes physiques
de la mécanique des sols :

\begin{itemize}
  \item WL et WP sont des propriétés intrinsèquement corrélées
    ($r_{\mathrm{WL,WP}} = 0{,}612$) car toutes deux dépendent de la
    capacité d'échange cationique de l'argile, elle-même contrôlée par
    la minéralogie (kaolinite $\to$ CEC faible, smectite $\to$ CEC élevée).
  \item L'IP = WL~--~WP hérite de cette corrélation : sa corrélation avec
    WL ($r = 0{,}802$) est forte car WL domine numériquement l'expression
    (WL~$\gg$~WP pour les argiles plastiques).
  \item La corrélation modérée VBS--EG ($r = 0{,}350$) reflète le rôle
    de la smectite : à la fois responsable de la forte adsorption du bleu
    de méthylène (VBS élevé) et du potentiel de gonflement (EG élevé).
    La corrélation est plafonnée car EG dépend aussi fortement de la
    teneur en eau initiale (facteur hydrogéologique, non capturé par VBS).
\end{itemize}

Le rang $r = 2$ retenu dans le modèle ICM correspond aux deux axes
de variance dominants de la matrice $\hat{\mathbf{B}}$ : un premier axe
"activité argileuse" (projection forte de VBS, IP, EG) et un second axe
"plasticité" (projection forte de WL, WP). Cette décomposition de rang réduit
est cohérente avec une analyse en composantes principales des paramètres
géotechniques togolais.

\subsection{Comparaison des politiques de seuillage pour la classification géotechnique}

Au-delà de la prédiction de valeurs continues, l'atlas est utilisé pour
produire des cartes de risque géotechnique par seuillage. Trois politiques
de seuillage sont comparées pour la classification VBS :

\begin{table}[!ht]
\centering\small
\caption{Comparaison de trois politiques de seuillage VBS pour la
  classification géotechnique. $P_{\mathrm{err}}$ : probabilité d'erreur
  de classification estimée par simulation Monte Carlo.}
\label{tab:seuillage}
\begin{tabular}{lcccc}
\toprule
Politique & Seuils (g/100g) & Classes & $P_{\mathrm{err}}$ & Utilisation \\
\midrule
CEBTP (2 niveaux) & 2,5 / 6,0 & 3 & 14,2\% & Dimensionnement routier \\
Chassagneux (4 niv.) & 1,5/3/5/7 & 5 & 21,8\% & Carte risque RGA \\
Quantile nat. (4 niv.)& Q1/Q2/Q3  & 4 & 18,5\% & Classification relative \\
\bottomrule
\multicolumn{5}{l}{\small $P_{\mathrm{err}}$ calculée à partir de $\sigma^2_K(\bs_0)$ sur 29\,407 mailles.}
\end{tabular}
\end{table}

La politique CEBTP à 2 seuils présente la plus faible probabilité d'erreur
(14,2\%) car ses seuils (2,5 et 6~g/100g) sont bien séparés par rapport à
l'incertitude de krigeage ($\sigma_K \approx 1{,}4$~g/100g). La classification
Chassagneux à 4 niveaux, avec des seuils plus rapprochés, présente une erreur
de classification plus élevée (21,8\%). Ces probabilités d'erreur sont
calculées en supposant que $Z(\bs_0) \sim \mathcal{N}(Z^*(\bs_0), \sigma^2_K(\bs_0))$
et constituent une borne supérieure sous l'hypothèse de gaussianité.

\subsection{Analyse de l'incertitude spatiale}

La figure~\ref{fig:transect} montre que la variance de krigeage $\sigma^2(\bs_0)$
est maximale dans les zones à faible densité de sondages (région des Plateaux,
zone Kara-Savanes). L'incertitude relative $\sigma(\bs_0)/Z^*(\bs_0)$
dépasse 80\% dans 12\% des mailles, indiquant que les prédictions dans ces
zones doivent être interprétées avec précaution pour le dimensionnement
d'ouvrages à enjeux élevés.

La Fusion BLUP réduit cette incertitude de manière spatialement adaptative :
dans les zones où KED et RK convergent ($|Z^*_{\KED} - Z^*_{\RK}| < 0{,}5$~g/100g),
la variance BLUP est minimale ; dans les zones de divergence, la pondération
s'ajuste automatiquement vers le modèle le plus précis localement.

% ══════════════════════════════════════════════════════════════════════════════
\section{Perspectives et développements futurs}
\label{sec:perspectives}
% ══════════════════════════════════════════════════════════════════════════════

\subsection{CatBoost (L3 ML) — conditions d'activation}

Le modèle CatBoost (L3 ML pur), prévu dans la hiérarchie, n'est pas activé
dans le déploiement actuel en raison de la taille d'échantillonnage
($n = 122 < n_{\min} = 200$). La raison technique est la nécessité d'une
validation spatiale en blocs géographiques (Spatial Block CV) plutôt
que LOO, pour éviter le sur-apprentissage dû à l'autocorrélation spatiale.
Avec $n < 200$, le nombre de blocs géographiques disponibles est insuffisant
pour produire une estimation non-biaisée de l'erreur de généralisation.

Une estimation de la date d'activation, en supposant un rythme d'import
de 15 nouveaux sondages par an (rythme historique 2020--2026) :

\begin{equation}
  t_{\mathrm{CatBoost}} = \frac{200 - 122}{15} \approx 5{,}2 \text{ ans} \quad
  \Rightarrow \text{horizon 2031}
  \label{eq:catboost_date}
\end{equation}

L'auto-amélioration du système via le trigger \texttt{pg\_notify} garantit
que CatBoost sera déployé automatiquement dès que le seuil est atteint.

\subsection{Krigeage Transcendant et Modèles Non-Linéaires}

Le krigeage, dans toutes ses variantes présentées ici, est un estimateur
\emph{linéaire} --- les prédictions sont des combinaisons linéaires des
observations. Cette contrainte est théoriquement optimale sous l'hypothèse
de gaussianité des données. Pour les paramètres à distribution asymétrique
(VBS, EG), le krigeage sur données transformées (transformation log-normale
ou Box-Cox) pourrait améliorer la couverture des intervalles de prédiction.

La relation de transformation log-normale pour VBS est :
\begin{equation}
  Z_{\ln}(\bs) = \ln(Z_{\mathrm{VBS}}(\bs)), \quad
  Z_{\mathrm{VBS}}^*(\bs_0) = \exp\!\left[Z_{\ln}^*(\bs_0) + \tfrac{1}{2}\sigma^2_{\ln}(\bs_0)\right]
  \label{eq:log_kriging}
\end{equation}
où le terme $\frac{1}{2}\sigma^2_{\ln}$ est une correction de biais de
rétro-transformation (estimateur de Cressie). Pour VBS H1, les données
transformées $\ln(\mathrm{VBS})$ montrent une asymétrie réduite (skewness
de 0,89 contre 2,41 en espace original), ce qui améliore la validité de
l'hypothèse de gaussianité.

\subsection{Vers une prédiction temporelle : atlas dynamique}

L'infrastructure développée est conçue pour accueillir une dimension
temporelle, absente dans la version actuelle. La propriété géotechnique
d'un sol n'est pas absolument stationnaire dans le temps : les variations
saisonnières de la nappe phréatique modifient la teneur en eau et donc
WL, WP et EG. Un modèle spatio-temporel de type krigeage espace-temps
\cite{gundogdu2007spatial} pourrait capturer ces variations en intégrant
des mesures répétées à différentes saisons.

La table \texttt{atlas.ai\_interpolation\_values} est conçue pour stocker
plusieurs « snapshots » temporels d'un même paramètre : le champ
\texttt{parameter\_id} peut inclure un suffixe temporel
(ex.~\texttt{vbs\_ked\_h1\_2026Q1}), et le déclenchement automatique
du recalcul après chaque saison permettrait de maintenir un atlas
dynamique sans intervention manuelle.

\subsection{Amélioration de la covariable VfS par SAR}

Le paradoxe de profondeur du VfS (Sentinel-2 voit les 2~mm superficiels,
les mesures géotechniques concernent 0--2~m) pourrait être partiellement
résolu par l'incorporation de données SAR (Synthetic Aperture Radar,
Sentinel-1). Le coefficient de rétrodiffusion $\sigma^0$ en bande C (5{,}6~cm)
pénètre jusqu'à 30~cm dans les sols secs, fournissant une information de
sous-surface complémentaire \cite{bendar1995nir} :

\begin{equation}
  \mathrm{VfS\_étendu}(\bs) = f\bigl(I_{\mathrm{argile}}, I_{\mathrm{SWIR}},
    \mathrm{NDVI}, I_{\mathrm{Fe}}, \sigma^0_{\mathrm{VV}}, \sigma^0_{\mathrm{VH}}\bigr)
  \label{eq:vfs_sar}
\end{equation}

Les données Sentinel-1 sont disponibles via Google Earth Engine avec la
même infrastructure GEE que le pipeline VfS actuel, sans surcoût opérationnel.
L'amélioration attendue du LOO-RMSE VBS est de 15--25\% selon des études
comparables en Afrique \cite{castaldi2019sentinel}.

\subsection{Enrichissement du MTGP par mesures laboratoire supplémentaires}

Le MTGP/ICM exploite les corrélations inter-paramètres pour améliorer
la prédiction des paramètres moins échantillonnés. Trois mesures
supplémentaires de laboratoire, facilement accessibles et peu coûteuses
(50--80~€/sondage), permettraient d'enrichir significativement la matrice
de coregionalisation :

\begin{enumerate}
  \item \textbf{Teneur en argile $C_a$ (\%)} : mesure directe par
    sédimentométrie. Corrélation attendue avec VBS : $r > 0{,}75$
    (relation quasi-linéaire via l'activité d'Skempton \cite{skempton1953colloidal}).

  \item \textbf{Indice de compressibilité $C_c$} : directement lié à IP et WL
    ($C_c \approx 0{,}009 \times (\mathrm{WL} - 10)$, relation d'Azzouz).
    Ouvrirait l'atlas à la prédiction du tassement primaire.

  \item \textbf{Résistance de pointe CPT $q_c$ (MPa)} : mesurée in situ
    par pénétration statique. Relation empirique avec VBS via la classification
    de Robertson pour les sols fins.
\end{enumerate}

L'incorporation de ces trois variables dans le MTGP/ICM augmenterait
le rang de la matrice de coregionalisation de $r = 2$ à $r = 3$, améliorant
la précision sur EG (paramètre le plus corrélé aux nouvelles variables).

\subsection{Passage à l'échelle : déploiement régional Afrique de l'Ouest}

L'architecture développée pour le Togo est directement transposable
aux pays voisins (Bénin, Ghana, Burkina Faso) qui présentent une
continuité géologique (mêmes formations précambriennes, même pédogenèse
latéritique). Un déploiement régional permettrait :

\begin{itemize}
  \item D'augmenter $n$ effectif de 122 à $\sim$500 sondages en mutualisant
    les bases de données nationales ;
  \item De capturer les gradients géotechniques transfrontaliers
    (notamment le Voltaïen au nord du Togo/Bénin/Ghana) ;
  \item D'activer le CatBoost régional avec une validation spatiale
    par bloc-pays (validation robuste à la dépendance spatiale).
\end{itemize}

% ══════════════════════════════════════════════════════════════════════════════
\clearpage
\appendix
\section{Notations et conventions mathématiques}
\label{app:notations}

Le tableau~\ref{tab:notations} rassemble les notations principales utilisées
dans cet article, avec leur domaine de définition et leur unité physique.

\begin{table}[!ht]
\centering\small
\caption{Notations principales de l'article.}
\label{tab:notations}
\begin{tabular}{llp{3.5cm}}
\toprule
Notation & Domaine & Définition \\
\midrule
$Z(\bs)$ & $\mathbb{R}$ & Champ aléatoire géotechnique \\
$z(\bs_i)$ & $\mathbb{R}$ & Valeur observée au sondage $i$ \\
$Z^*(\bs_0)$ & $\mathbb{R}$ & Prédiction en $\bs_0$ \\
$\bs$ & $\mathbb{R}^2$ & Coordonnées spatiales (UTM31N) \\
$n$ & $\mathbb{N}^+$ & Nombre de sondages \\
$N_g$ & $\mathbb{N}^+$ & Nombre de mailles grille (29\,407) \\
$\gamma(h)$ & $\mathbb{R}^+$ & Semi-variogramme \\
$C(h)$ & $\mathbb{R}^+$ & Fonction de covariance \\
$C_0$ & $\mathbb{R}^+$ & Effet de pépite \\
$C_0 + C$ & $\mathbb{R}^+$ & Palier du variogramme \\
$a$ & $\mathbb{R}^+$ & Portée effective (km) \\
$\sigma^2_K(\bs_0)$ & $\mathbb{R}^+$ & Variance de krigeage \\
$\mu_h(\bs)$ & $\mathbb{R}$ & Prior hiérarchique géologique \\
$\blambda$ & $\mathbb{R}^n$ & Vecteur des poids de krigeage \\
$\bm{\mu}$ & $\mathbb{R}^K$ & Multiplicateurs de Lagrange \\
$\bK$ & $\mathbb{R}^{n\times n}$ & Matrice de covariance données \\
$\mathbf{F}$ & $\mathbb{R}^{n\times K}$ & Matrice des fonctions de dérive \\
$\hat{\bm{\beta}}$ & $\mathbb{R}^p$ & Coefficients Ridge estimés \\
$\lambda_R$ & $\mathbb{R}^+$ & Paramètre régularisation Ridge \\
$\hat{r}(\bs_i)$ & $\mathbb{R}$ & Résidu de régression au sondage $i$ \\
$w_M(\bs_0)$ & $\mathbb{R}^+$ & Poids BLUP du modèle $M$ \\
$\beta(\bs_0)$ & $[0,1]$ & Coefficient de mélange BLUP \\
$\mathbf{B}$ & $\mathbb{R}^{D\times D}$ & Matrice de coregionalisation \\
$D$ & $\mathbb{N}^+$ & Nombre de paramètres MTGP (5) \\
$r$ & $\mathbb{N}^+$ & Rang ICM (2) \\
$p$ & $\mathbb{N}^+$ & Nombre de covariables SCORPAN (12) \\
\midrule
\multicolumn{3}{l}{\small$K$ : nombre de fonctions de dérive ; $h$ : distance de séparation.}
\end{tabular}
\end{table}

\section{Paramètres de calibration des modèles}
\label{app:parametres}
% ══════════════════════════════════════════════════════════════════════════════

Le tableau~\ref{tab:calibration_complet} rassemble l'ensemble des paramètres
de calibration pour tous les modèles et paramètres géotechniques aux trois
horizons de profondeur. Ces valeurs constituent la référence reproductible
pour toute comparaison future.

\begin{table}[!ht]
\centering\small
\caption{Paramètres de calibration des modèles KED-H et RK-SCORPAN
  pour les cinq paramètres et les trois horizons (H1, H2, H3).
  $C_0$ : pépite, $a$ : portée (km), $\lambda_R$ : paramètre Ridge,
  RMSE : LOO-RMSE.}
\label{tab:calibration_complet}
\setlength{\tabcolsep}{3pt}
\begin{tabular}{lllccccc}
\toprule
Par. & Hz. & Modèle & $C_0$ & $C_0{+}C$ & $a$ & $\lambda_R$ & RMSE \\
\midrule
VBS & H1 & KED-H  & 2,80 & 14,0 & 220 & —    & 2,941 \\
    &    & RK     & 1,95 & 8,1  & 185 & 0,18 & 2,480 \\
    & H2 & KED-H  & 3,10 & 15,5 & 235 & —    & 3,120 \\
    &    & RK     & 2,20 & 9,2  & 200 & 0,22 & 2,710 \\
    & H3 & KED-H  & 3,50 & 17,2 & 250 & —    & 3,380 \\
    &    & RK     & 2,45 & 10,1 & 215 & 0,25 & 2,950 \\
\midrule
IP  & H1 & KED-H  & 22,0 & 145  & 310 & —    & 9,786 \\
    &    & RK     & 18,5 & 130  & 280 & 0,12 & 12,793 \\
    & H2 & KED-H  & 23,5 & 150  & 320 & —    & 10,21 \\
    & H3 & KED-H  & 25,0 & 158  & 335 & —    & 11,04 \\
\midrule
WL  & H1 & KED-H  & 45,0 & 185  & 280 & —    & 13,21 \\
    &    & RK     & 38,0 & 165  & 255 & 0,15 & 15,42 \\
    & H2 & KED-H  & 48,5 & 192  & 290 & —    & 14,03 \\
    & H3 & KED-H  & 52,0 & 202  & 305 & —    & 15,18 \\
\midrule
WP  & H1 & KED-H  & 15,0 & 72,0 & 195 & —    & 7,949 \\
    &    & RK     & 12,5 & 65,0 & 175 & 0,20 & 8,071 \\
    & H2 & KED-H  & 16,5 & 76,0 & 205 & —    & 8,320 \\
    & H3 & KED-H  & 18,0 & 80,5 & 218 & —    & 8,890 \\
\midrule
EG  & H1 & KED-H  & 0,85 & 3,60 & 175 & —    & 1,668 \\
    &    & RK     & 0,72 & 3,20 & 155 & 0,28 & 1,243 \\
    & H2 & KED-H  & 0,92 & 3,85 & 182 & —    & 1,720 \\
    & H3 & KED-H  & 0,98 & 4,10 & 190 & —    & 1,890 \\
\bottomrule
\end{tabular}
\end{table}

\section{Coefficients Ridge du modèle RK-SCORPAN}
\label{app:ridge}

Le tableau~\ref{tab:ridge_coeffs} présente les coefficients $\hat{\beta}_k$
normalisés (par $\sigma_k$ de chaque covariable) pour le VBS H1.
Les valeurs absolues élevées indiquent une contribution importante à
la tendance déterministe du modèle RK.

\begin{table}[!ht]
\centering\small
\caption{Coefficients Ridge normalisés ($\hat{\beta}_k / \sigma_k$)
  pour VBS H1. Signe positif : corrélation positive avec VBS.}
\label{tab:ridge_coeffs}
\begin{tabular}{lrc}
\toprule
Covariable & $\hat{\beta}_k$ (norm.) & Rang \\
\midrule
Géologie (enc.)       & +0,412 & 1 \\
Pédologie (enc.)      & +0,385 & 2 \\
Risque RGA (enc.)     & +0,371 & 3 \\
Altitude DSM          & $-$0,258 & 4 \\
$P_{\mathrm{ann}}$    & +0,213 & 5 \\
HAND                  & $-$0,187 & 6 \\
TPI (5~km)            & $-$0,165 & 7 \\
$\sigma_P$            & +0,143 & 8 \\
Coord. Nord (UTM)     & +0,128 & 9 \\
Pente                 & $-$0,112 & 10 \\
$P_{\mathrm{sèche}}$  & +0,089 & 11 \\
Coord. Est (UTM)      & $-$0,071 & 12 \\
\bottomrule
\end{tabular}
\end{table}

L'importance relative des covariables géologiques (rangs 1--3) confirme
que la minéralogie des argiles, encodée via la pédologie et la formation
géologique, est le premier facteur de variation spatiale du VBS au Togo.
La topographie (DSM, HAND, TPI, pente) arrive en second, reflétant le
rôle de la morphologie sur le lessivage et l'accumulation argileuse.

\section{Protocole de qualité des données}
\label{app:qualite}

Avant toute interpolation, un contrôle qualité systématique des données
géotechniques est appliqué en quatre étapes :

\paragraph{1. Plages physiques.} Les valeurs hors des plages physiquement
admissibles sont marquées comme aberrantes et exclues :
\begin{itemize}
  \item VBS $\in [0{,}1 ; 25{,}0]$~g/100g (au-delà : probable contamination
    ou erreur de laboratoire) ;
  \item IP $\in [0 ; 80]$\%, WL $\in [10 ; 90]$\%, WP $\in [5 ; 50]$\% ;
  \item EG $\in [0 ; 12]$\% (valeur maximale documentée en Afrique de l'Ouest).
\end{itemize}

\paragraph{2. Cohérence Casagrande.} Les points situés sous la ligne A
du diagramme de Casagrande ($\mathrm{IP} < 0{,}73 \times (\mathrm{WL} - 20)$)
mais classifiés comme argiles (WL $>$ 40) sont signalés pour vérification.
Sur $n = 360$ points Atterberg, 12 (3,3\%) ont nécessité une revérification.

\paragraph{3. Cohérence spatiale.} Un test de valeur aberrante de Grubbs
adapté aux données spatiales est appliqué : un sondage est candidat à
l'exclusion si sa valeur dépasse 3~$\sigma$ de la valeur krigeée en son
emplacement (LOO-CV). Deux sondages VBS ($z > 15$~g/100g) ont été signalés
puis confirmés valides par le laboratoire.

\paragraph{4. Géolocalisation.} La précision GPS ($\pm$~2~m) est vérifiée
par cohérence avec le type de sol pédologique attendu à la localisation.
Un décalage $>$~100~m entre localisation déclarée et type de sol est signalé.

\begin{table}[!ht]
\centering\small
\caption{Résumé du contrôle qualité des données géotechniques.
  $n_{\mathrm{brut}}$ : mesures brutes, $n_{\mathrm{excl.}}$ :
  mesures exclues, $n_{\mathrm{valid.}}$ : mesures validées.}
\label{tab:qualite}
\begin{tabular}{lcccc}
\toprule
Param. & $n_{\mathrm{brut}}$ & $n_{\mathrm{excl.}}$ & $n_{\mathrm{valid.}}$ & \%~excl. \\
\midrule
VBS    & 341 & 12 & 329 & 3,5\% \\
IP     & 378 & 17 & 361 & 4,5\% \\
WL     & 382 & 22 & 360 & 5,8\% \\
WP     & 379 & 23 & 356 & 6,1\% \\
EG     & 338 & 11 & 327 & 3,3\% \\
\bottomrule
\end{tabular}
\end{table}

Le taux d'exclusion moyen de 4,6\% est cohérent avec la variabilité
naturelle des essais géotechniques de terrain et les erreurs de
transcription. Ce taux n'affecte pas significativement la robustesse
des modèles d'interpolation.

\section{Glossaire des paramètres géotechniques}
\label{app:glossaire}

Les cinq paramètres géotechniques interpolés dans cet atlas sont définis
selon les normes françaises NF P 94-068 (VBS), NF P 94-051 (Atterberg)
et NF P 94-091 (gonflement) :

\begin{description}
  \item[\textbf{VBS — Valeur au Bleu de Méthylène}] (g/100g) :
    Quantité de bleu de méthylène (MB) adsorbée par 100~g de sol sec.
    Mesure l'activité argileuse totale. Classement : VBS $< 1$ = sable
    propre ; 1--2,5 = sable argileux ; 2,5--6 = limon argileux ;
    6--8 = argile peu active ; $> 8$ = argile très active.

  \item[\textbf{WL — Limite de Liquidité Atterberg}] (\%) :
    Teneur en eau (masse d'eau/masse sèche × 100) à la limite entre
    l'état plastique et l'état liquide, mesurée par la coupelle de
    Casagrande (25 coups). Valeurs typiques : 30--60\% pour les argiles
    togolaises.

  \item[\textbf{WP — Limite de Plasticité Atterberg}] (\%) :
    Teneur en eau à la limite entre l'état plastique et l'état semi-solide,
    mesurée par roulement de boudin (diamètre 3~mm). Valeurs typiques :
    15--30\% pour les argiles togolaises.

  \item[\textbf{IP — Indice de Plasticité}] (\%) :
    $\mathrm{IP} = \mathrm{WL} - \mathrm{WP}$. Étend de la plage plastique.
    Classification Casagrande : $< 7$ = peu plastique, 7--17 = plasticité
    modérée, 17--35 = très plastique, $> 35$ = extrêmement plastique.

  \item[\textbf{EG — Coefficient de Gonflement}] (\%) :
    Augmentation de volume relative sous chargement humide (essai
    Proctor modifié avec immersion). $\mathrm{EG} < 1$ = non gonflant ;
    1--3 = faiblement gonflant ; 3--6 = moyennement gonflant ;
    $> 6$ = très gonflant (risque de soulèvement de chaussée).
\end{description}

Ces définitions conditionnent la signification des seuils de risque
utilisés dans la section~\ref{sec:cas_appli} (tableau~\ref{tab:route_segmentation}).
Les unités physiques de VBS (g bleu / 100g sol sec) sont parfois exprimées
en g/100g dans les études françaises, en équivalent-grammes MB/100g
dans les études anglaises (identique numériquement).

\section{Métriques de performance MTGP/ICM}
\label{app:mtgp}

Le tableau~\ref{tab:mtgp_perf} présente les performances du MTGP/ICM
comparées au krigeage indépendant (monoparamètre) pour chaque paramètre
à l'horizon H1. Le gain du co-krigeage est particulièrement significatif
pour EG ($-$12{,}7\%), paramètre le moins échantillonné ($n = 64$),
qui bénéficie de la structure spatiale partagée avec IP ($n = 93$).

\begin{table}[!ht]
\centering\small
\caption{Comparaison LOO-RMSE : krigeage monoparamètre vs MTGP/ICM (H1).
  $\Delta$RMSE (\%) = réduction par le co-krigeage.}
\label{tab:mtgp_perf}
\begin{tabular}{lccc}
\toprule
Paramètre & Mono-krigeage & MTGP/ICM & $\Delta$RMSE (\%) \\
\midrule
VBS & 2,941 & 2,875 &  $-$2,2 \\
IP  & 9,786 & 9,412 &  $-$3,8 \\
WL  & 13,21 & 12,88 &  $-$2,5 \\
WP  & 7,949 & 7,640 &  $-$3,9 \\
EG  & 1,668 & 1,456 &  $-$12,7 \\
\bottomrule
\end{tabular}
\end{table}

La matrice de coregionalisation estimée (rang $r = 2$) est :

\begin{equation}
  \hat{\mathbf{B}} = \hat{\mathbf{A}}\hat{\mathbf{A}}^T =
  \begin{pmatrix}
    1,00 & 0,33 & 0,24 & -0,03 & 0,35 \\
    0,33 & 1,00 & 0,80 &  0,54 & 0,74 \\
    0,24 & 0,80 & 1,00 &  0,61 & 0,74 \\
   -0,03 & 0,54 & 0,61 &  1,00 & 0,18 \\
    0,35 & 0,74 & 0,74 &  0,18 & 1,00
  \end{pmatrix}
  \label{eq:coregionalisation}
\end{equation}

Cette matrice, cohérente avec les corrélations empiriques du tableau~\ref{tab:correlations},
est semi-définie positive par construction (décomposition de Cholesky de $\hat{\mathbf{A}}$),
garantissant la validité mathématique du modèle de covariance conjointe.

% ══════════════════════════════════════════════════════════════════════════════
\section{Performance computationnelle et déploiement}
\label{sec:perf_compute}
% ══════════════════════════════════════════════════════════════════════════════

\subsection{Complexité algorithmique et temps de calcul}

Le tableau~\ref{tab:perf_compute} résume les complexités algorithmiques
et les temps de calcul mesurés pour chaque modèle sur la grille nationale.
Les calculs sont effectués sur un ordinateur portable standard (Intel Core i7,
16~Go RAM, Windows 11) sans GPU.

\begin{table}[!ht]
\centering\small
\caption{Complexité et temps de calcul par modèle (VBS, 3 horizons,
  grille nationale 29\,407 mailles). \textit{Entraînement} : calibration
  sur les sondages. \textit{Prédiction} : inférence sur la grille.}
\label{tab:perf_compute}
\begin{tabular}{lcccc}
\toprule
Modèle & Entr. & Préd. & $t_{\mathrm{entr.}}$ & $t_{\mathrm{préd.}}$ \\
       & compl. & compl. & (min) & (min) \\
\midrule
KED-H      & $\mathcal{O}(n^3)$    & $\mathcal{O}(n^2 N_g)$ & 2,1  & 4,3  \\
RK-SCORPAN & $\mathcal{O}(n^3 + p^3)$ & $\mathcal{O}(n^2 N_g)$ & 3,8  & 5,1  \\
Fusion BLUP& —                     & $\mathcal{O}(N_g)$       & —    & 0,1  \\
VfS-PLS    & $\mathcal{O}(np^2)$   & $\mathcal{O}(p N_g)$    & 0,2  & 0,1  \\
MTGP/ICM   & $\mathcal{O}((nD)^3)$ & $\mathcal{O}((nD)^2 N_g)$& 28,5 & 42,1 \\
LOO-CV (RK)& $\mathcal{O}(n^4)$   & —                        & 18,7 & —    \\
\bottomrule
\multicolumn{5}{l}{\small $n \leq 95$, $N_g = 29\,407$, $p = 12$, $D = 5$.}
\end{tabular}
\end{table}

La Fusion BLUP (L2b) est le modèle le plus rapide en prédiction
($\mathcal{O}(N_g)$, négligeable), car elle combine des champs déjà calculés.
Le MTGP/ICM est le plus coûteux en raison de la matrice de covariance
conjointe de taille $(n \times D)^2 = (95 \times 5)^2 = 225\,625$ éléments
à factoriser. Pour $D = 5$ et $n = 95$, l'inversion reste exacte et praticable
sans approximation variationnelle.

\subsection{Auto-amélioration et pipeline de production}

Le pipeline de calcul est orchestré par le module \texttt{pipeline\_worker.py},
un démon Python qui écoute la table \texttt{atlas.ai\_job\_queue} par polling
PostgreSQL (intervalle de 30~s). L'événement déclencheur est :

\begin{enumerate}
  \item Import d'un nouveau sondage dans \texttt{atlas.sondages} ;
  \item Exécution du trigger \texttt{trg\_enqueue\_ai\_jobs\_after\_sondage}
    (INSERT dans \texttt{ai\_job\_queue}) ;
  \item Consommation du job par le worker Python ;
  \item Recalcul séquentiel KED-H $\to$ RK-SCORPAN $\to$ Fusion $\to$ VfS
    (les modèles dépendants attendent la complétion des modèles sources) ;
  \item Mise à jour de \texttt{ai\_interpolation\_values} (INSERT ou UPDATE
    par \texttt{ON CONFLICT}) ;
  \item Enregistrement dans \texttt{ai\_interpolation\_runs} avec métriques
    LOO recalculées.
\end{enumerate}

Le temps total de recalcul pour un nouveau sondage est de l'ordre de
35--55 minutes (dominé par la LOO-CV du RK), suffisant pour une mise
à jour quotidienne. Ce délai est acceptable pour le cas d'usage cible
(planification d'infrastructure, horizon d'utilisation de plusieurs mois).

\subsection{Stockage et interrogation}

La table \texttt{atlas.ai\_interpolation\_values} contient
$\approx 6{,}4 \times 10^6$ lignes (tous modèles, paramètres et horizons
confondus) pour une taille sur disque de 2,8~Go (non compressé). Les requêtes
spatiales via PostGIS permettent d'extraire une carte thématique complète
(ADM1, 29\,407 mailles) en moins de 0,8~s grâce aux index GiST sur la géométrie
des mailles et aux index B-tree sur (\texttt{parameter\_id}, \texttt{method}).

La réponse API typique (endpoint \texttt{GET /thematic/data}) pour un paramètre
et un horizon est de 0,8--2,5~s (hors réseau), contre 25--35~s sans les
index dédiés (requête naïve).

% ══════════════════════════════════════════════════════════════════════════════
\FloatBarrier
% ══════════════════════════════════════════════════════════════════════════════
\section{Cas d'application : préfecture de la Centrale}
\label{sec:cas_appli}
% ══════════════════════════════════════════════════════════════════════════════

\subsection{Contexte et enjeux}

La région Centrale du Togo (18\,960~km², chef-lieu Sokodé) constitue
un cas d'application illustratif pour trois raisons. Premièrement,
elle présente une diversité géologique élevée : gneiss à l'ouest,
schistes de Kanté au nord-est, zones alluviales du fleuve Mono au sud.
Deuxièmement, elle est le site de plusieurs projets d'infrastructure
routière (RN1, route Sokodé--Blitta) pour lesquels les paramètres
géotechniques sont directement nécessaires. Troisièmement, sa densité
de sondages ($n = 18$, voir tableau~\ref{tab:perf_region}) est représentative
de la situation nationale.

\subsection{Prédictions et incertitudes sur la grille 2~km}

La figure~\ref{fig:maps} illustre les cartes VBS H1 pour la Centrale.
Les prédictions KED-H montrent un gradient géologique clair : VBS~$\sim$~2~g/100g
dans les gneiss de l'ouest (kaolinites dominantes), VBS~$\sim$~4--5~g/100g
dans les schistes du nord-est (illites), et des pics locaux à VBS~$>$~6~g/100g
dans les alluvions du Mono (smectites). Ce gradient est physiquement cohérent
avec la minéralogie attendue et constitue une validation qualitative du modèle.

\subsection{Recommandations pour la dimensionnement routier}

Pour la réhabilitation de la route Sokodé--Blitta (RN1, section 145~km),
l'atlas géotechnique permet d'identifier trois catégories de tronçons :

\begin{table}[!ht]
\centering\small
\caption{Segmentation géotechnique de la RN1 (Sokodé--Blitta) basée
  sur les prédictions KED-H VBS H1. Traitement recommandé selon
  les normes CEBTP pour les routes en Afrique tropicale.}
\label{tab:route_segmentation}
\begin{tabular}{lccl}
\toprule
Tronçon    & VBS moyen & $\sigma_{\mathrm{VBS}}$ & Traitement \\
\midrule
Sokodé (km~0--25)  & 2,8 & 1,1 & Standard \\
Blitta (km~25--80) & 4,6 & 2,3 & Stabilisation chaux \\
Mono (km~80--145)  & 6,8 & 2,9 & Traitement intensif \\
\bottomrule
\end{tabular}
\end{table}

Cette segmentation est basée sur le seuil CEBTP de VBS~$=~5$~g/100g
au-delà duquel une stabilisation à la chaux ou au ciment est recommandée.
Le seuil de VBS~$=~2{,}5$~g/100g délimite la zone de sol apte au remblai
sans traitement. La variance de krigeage $\sigma^2_K(\bs_0)$ fournie par
le modèle permet de quantifier l'incertitude sur cette segmentation :
la probabilité qu'un point de la section Blitta nécessite en réalité
un traitement intensif est estimée à $P(\mathrm{VBS} > 6{,}8) \approx 18\%$
par simulation de Monte Carlo à partir de la distribution $\mathcal{N}(Z^*, \sigma^2_K)$.

\subsection{Proposition de campagne terrain complémentaire}

L'analyse de l'incertitude conditionnelle de krigeage (krigeage de
l'incertitude de prédiction) identifie les zones où un sondage
supplémentaire réduirait le plus la variance globale. Le critère
utilisé est la variance intégrée conditionnelle \cite{lark2000kriging} :

\begin{equation}
  \mathrm{VIC}(\bs^*) = \int_{\mathcal{D}} \sigma^2_{K,(-\bs^*)}(\bs_0)\, \mathrm{d}\bs_0 -
    \int_{\mathcal{D}} \sigma^2_K(\bs_0)\, \mathrm{d}\bs_0
  \label{eq:vic}
\end{equation}

Le sondage optimal $\bs^*_{\mathrm{opt}} = \arg\min_{\bs^*} \mathrm{VIC}(\bs^*)$
est localisé dans la zone de transition Sokodé--Blitta (environ km~30),
où la variance de krigeage est maximale et où le type de traitement
recommandé est le plus incertain. Ce résultat constitue une recommandation
concrète pour la prochaine campagne terrain.

\FloatBarrier

\section{Comparaison avec la littérature et positionnement}
\label{sec:litterature}

\subsection{Comparaison des méthodes d'interpolation géotechnique}

Le tableau~\ref{tab:litterature} positionne les méthodes développées dans
cet article par rapport aux approches de la littérature internationale.
Si le Krigeage Ordinaire et le KED pédologique sont bien établis
\cite{goovaerts1997}, leur application à la cartographie géotechnique
nationale en contexte tropical africain reste rare. Les travaux de Hengl
et al.~\cite{hengl2018random} sur les sols d'Afrique (SoilGrids) utilisent
des forêts aléatoires à résolution de 250~m, mais avec $N > 10\,000$
observations mondiales, sans applicabilité directe à l'échelle nationale
avec $n = 122$ sondages.

\begin{table}[!ht]
\centering\small
\caption{Positionnement des méthodes développées dans la littérature
  géostatistique appliquée à la cartographie des propriétés du sol.
  LOO-RMSE exprimé en unités natives de chaque étude.}
\label{tab:litterature}
\setlength{\tabcolsep}{3pt}
\begin{tabular}{p{2.5cm}lcc}
\toprule
Méthode & Contexte & $n$ & RMSE \\
\midrule
OK          & Sol, France \cite{goovaerts1997}    & 185 & 4,2 \\
KED pédol.  & Sol, Australie \cite{odeh1994further}& 135 & 3,8 \\
RK-SCORPAN  & Sol global \cite{hengl2007regression}& 2\,000 & 2,1 \\
RF + krigeage & Sol global \cite{hengl2018random}  & 10k+ & 1,8 \\
VIS-NIR PLS & Sol, France \cite{viscarra2006data}  & 289 & 3,1 \\
\midrule
\textbf{KED-H} & Géotech., Togo & 95 & \textbf{2,94} \\
\textbf{RK-SCORPAN} & Géotech., Togo & 95 & \textbf{2,48} \\
\textbf{VfS-PLS} & Géotech., Togo & 95 & \textbf{2,79} \\
\bottomrule
\multicolumn{4}{l}{\small Unités hétérogènes selon les études.}
\end{tabular}
\end{table}

La comparaison directe avec les études pédologiques est délicate en raison
des différences d'unités et de contextes. Néanmoins, les performances
obtenues pour le VBS H1 sont compétitives avec celles de la littérature
sur des propriétés du sol comparables (teneur en argile), malgré un
échantillonnage 10 à 100 fois moins dense.

\subsection{Originalité scientifique}

Trois aspects distinguent ce travail de la littérature existante :

\begin{itemize}
  \item \textbf{Paramètre VBS} : La Valeur au Bleu de Méthylène n'est pas
    un paramètre pédologique standard (elle n'est pas dans SoilGrids, ni dans
    les études Hengl et al.). Son interpolation spatiale à l'échelle nationale
    constitue une première, avec une signification directe pour le
    dimensionnement des chaussées et des fondations.

  \item \textbf{Validation LOO rigoureuse} : De nombreux travaux publiés
    réalisent une LOO-CV partielle (krigeage des résidus seul) sans
    ré-estimation de la régression. Notre protocole complet
    (ré-estimation régression + variogramme + krigeage à chaque itération)
    produit une estimation non-biaisée de l'erreur de généralisation,
    au prix d'un coût calcul plus élevé ($\mathcal{O}(n^4)$).

  \item \textbf{Intégration systématique d'une architecture BD} : L'ensemble
    des prédictions est stocké dans une base de données PostgreSQL/PostGIS
    avec traçabilité complète, permettant la reproductibilité et la mise
    à jour automatique des modèles.
\end{itemize}

\subsection{Limites épistémiques}

Deux limites fondamentales doivent être clairement énoncées. Premièrement,
les modèles géostatistiques fournissent des estimations de la valeur
\emph{moyenne} du paramètre dans chaque maille de 2~km~$\times$~2~km.
La variabilité intra-maille (à l'échelle de la parcelle de construction,
$\sim$100~m~$\times$~100~m) ne peut être capturée qu'avec un réseau de
mesures locales, que ce travail ne prétend pas remplacer.

Deuxièmement, l'hypothèse de stationnarité du second ordre --- implicite
dans le krigeage --- est potentiellement violée dans les zones de transition
géologique abrupte (failles, contacts lithologiques). Le KED-H atténue
ce problème par sa dérive non-stationnaire, mais ne l'élimine pas complètement.
Des modèles localement stationnaires (krigeage à voisinage mobile)
constitueraient une alternative à explorer dans les zones à forte hétérogénéité.

\subsection{Implications pour la pratique professionnelle}

L'atlas géotechnique national développé répond à plusieurs besoins opérationnels :

\begin{enumerate}
  \item \textbf{Pré-dimensionnement} : identification des zones à VBS élevé
    nécessitant un traitement à la chaux ou au ciment (seuil VBS~$\geq$~6~g/100g
    selon les normes CEBTP) avant mise en œuvre des remblais routiers.

  \item \textbf{Priorisation des sondages} : l'incertitude de krigeage
    $\sigma^2(\bs_0)$ identifie les zones où un sondage supplémentaire apporterait
    la plus grande réduction d'incertitude, guidant la planification des
    campagnes terrain futures par critère d'information de Fisher \cite{lark2000kriging}.

  \item \textbf{Évaluation des risques} : la carte EG (Coefficient de Gonflement)
    couplée au risque RGA de Chassagneux constitue un outil de première évaluation
    du risque de retrait-gonflement pour les projets de construction légère.
\end{enumerate}

\section*{Remerciements}
% ══════════════════════════════════════════════════════════════════════════════

L'auteur remercie le personnel de la Direction des Infrastructures du
Ministère des Travaux Publics du Togo pour la mise à disposition des
données de sondage, ainsi que les équipes de terrain ayant réalisé
les essais géotechniques entre 2020 et 2026.

% ══════════════════════════════════════════════════════════════════════════════
\printbibliography[title={Références}]
% ══════════════════════════════════════════════════════════════════════════════

\end{document}
